
\input grafinp3
%\input grafinput8
\input psfig
%\eqnotracetrue

\input psfig

\def\toone{{t+1}}
\def\ttwo{{t+2}}
\def\tthree{{t+3}}
\def\Tone{{T+1}}
\def\TTT{{T-1}}
\def\rtr{{\rm tr}}
\def\tone{{t+1}}

%\showchaptIDtrue
%\def\@chaptID{8.}

%\hbox{}
\footnum=0

  \auth{Samuelson, Paul}%
\chapter{Overlapping Generations\label{ogmodels}}

This chapter describes the pure exchange overlapping generations
model of Paul Samuelson (1958).
 We begin with an
abstract presentation  that treats the overlapping
generations model as a special case of the chapter \use{recurge}
general equilibrium model with complete markets and all
trades occurring at time $0$.
There is a special  type of heterogeneity across
agents: each person 
cares about consumption only at  two adjacent dates, and the set of
individuals who care about consumption at a particular date
includes some who care about consumption one period earlier and others
who care about consumption one period later.   We shall study
how  this
special preference and demographic  pattern affects some
 outcomes of the chapter \use{recurge} model.

  Imagining
complete markets with time $0$ trading in
an overlapping generations model strains
credulity by imagining  that equilibrium price and quantity sequences
are set at time $0$,
before most of the people who are to execute  trades have been born.
  For that reason, most applied work with the overlapping generations
model assumes a sequential-trading arrangement, like
the sequential trade in Arrow securities described   in chapter
\use{recurge}.
In a  sequential-trading arrangement,  all trades are 
agreed to only by living agents.

Despite its implausibility, we use the abstraction of
complete markets with time 0 trading to help us  understand
properties of   overlapping generations models 
that include  possible  multiplicity of equilibria and Pareto inefficient
equilibrium allocations, outcomes that don't  happen in the 
chapter \use{recurge} general equilibrium model. Since we  assume that everyone can 
trade with anyone else at time 0, we'll discover that such  outcomes 
 don't  hinge on 
agents alive today not being able to trade with agents not yet alive.\NFootnote{In section \use{sec:time0trading}, under
time $0$ trading, we will describe  an alternative 
equilibrium concept that entails a clearinghouse.}

%Lars%
%Nevertheless, equilibrium quantities and intertemporal prices
%are  equivalent between these  two trading arrangements, so 
%insights discovered in one setting transfer to the other.


Later in this chapter,
  we use versions of the model with
sequential trading  to tell how the overlapping
generations model provides  a framework for thinking about equilibria
with government debt, valued fiat currency, intergenerational
transfers, and fiscal policy.

%Lars%
% in sections \use{sec:Samuelson101}, \use{sec:deffinance},
%\use{sec:MMgov}, and \use{sec:realbills}.


\section{Preferences and endowments}

Time is discrete, starts at $t=0$,
and  lasts forever, so $t=0, 1, 2, \ldots$. In each period, 
there is a large number $N$ of identical agents born who 
live for two periods. The assumption of a large cohort
$N$ justifies an  analysis of a competitive equilibrium
in which agents take equilibrium prices as given. 
For expositional simplicity
we assume that there exists  a representative agent in each generation.
There is an infinity of representative agents,
one for each generation,
named $i=-1, 0, 1, \ldots$.  We can also regard $i$ as 
agent $i$'s period of birth.

  There is a single good at each date. The good is not storable.  There is no uncertainty.
  Each agent has a
strictly concave, twice  continuously
differentiable, one-period utility function $u(c)$ that
is strictly increasing in consumption  $c$ of the one good.
  Agent $i$ orders consumption stream  
${\bf c}^i = \{c_t^i\}_{t=0}^\infty$ according to the
  special  utility functional
$$\EQNalign{U^i({\bf c}^i) & = u(c_i^i) + u(c_{i+1}^i), \quad i \geq 0,
 \EQN olg1;a \cr
    U^{-1}({\bf c}^{-1}) &= u(c_0^{-1}). \EQN olg1;b \cr} $$
As a superscript, $i$ denotes a person's birth date; as a subscript, it denotes the current date.  
Notice that agent $i \geq 0 $ only wants goods dated $i$ and $i+1$, while an agent born at $i=-1$, before the
model starts at date $i =0$, wants consumption only at time $0$.
Equations \Ep{olg1} tell us  that agents
%Lars%  $i$ lives during periods $i$ and $i+1$ and wants 
want to consume only when they are alive in periods $t\geq 0$.

Each consumer $i \geq -1$ has an endowment sequence 
${\bf y}^i = \{y_t^i\}_{t=0}^\infty$.

\medskip
\noindent{\sc Definitions:}   An {\it allocation} is
a list $\{{\bf c}^i\}_{i\geq-1}
=\{\{c_t^i\}_{t\geq0}\}_{i\geq-1}$ of each agent $i$'s
consumption at every date $t$. 
A {\it feasible allocation} satisfies
$$\sum_i c_t^i \leq \sum_i y_t^i  \EQN olg_feasible $$
for all $t\geq 0$.
\medskip
\noindent 
In most of our analysis, we will assume that consumers 
are endowed with goods only when they are alive. 
That is, each consumer $i \geq 0$ has an endowment
sequence ${\bf y}^i$ satisfying
$y_i^i \geq 0$, $y^i_{i+1} \geq 0$, 
$y_t^i = 0 \ \forall t \neq i \ {\rm or}
\ i+1$; whereas the initial old consumer $i=-1$ has
a single endowment $y^{-1}_0 \geq 0$. 
Evidently an allocation is then feasible if for 
all $t \geq 0$,
$$ c_t^t + c_t^{t-1} \leq y_t^t + y_t^{t-1}. \EQN olg3 
$$
Furthermore, we will mostly focus on economies with 
stationary endowments in which  $y^i_i=y_y$ and 
$y^{i-1}_i=y_o$
for all $i\geq 0$. Here  subscript $y$ denotes
young and $o$ denotes old. 


\section{Dynamic (in)efficiency}\label{sec:dynineff}%
The Pareto inefficient equilibria encountered  in this
chapter reflect intertemporal inefficiencies.
Such equilibria  are commonly called 
{\it dynamically inefficient}. Under the preference
specification $U^i$ in \Ep{olg1} and stationary
endowments, $(y^i_i, y^{i-1}_i) = (y_y,y_o)$ for
all $i\geq 0$ and $y^i_t = 0$ otherwise,
multiple equilibria and dynamically inefficient 
equilibria  both exist 
if and only if $y_o<y_y$. If $y_o \geq y_y$, 
there is a unique equilibrium that is also Pareto 
efficient, as we found more generally for the  chapter \use{recurge} general 
equilibrium model.


\section{Time $0$ trading}\label{sec:time0trading}%
As in chapter \use{recurge}, we begin by analyzing an 
economy in which all trades take place at time $0$. 
But instead of a ``Walrasian auctioneer'' posting 
equilibrium prices we assume that there is a
``clearinghouse'' that posts prices, monitors demands 
and supplies for goods in all periods,
and verifies that each person's single time $0$ budget 
constraint is satisfied.
An equilibrium price vector assures that for
periods $t\geq 1$ aggregate endowments in that period 
are consumed by living agents,  but it can turn out 
that in period $0$ the 
aggregate endowment exceeds aggregate 
consumption which would mean that the clearinghouse
ends up with some time $0$ good not allocated.
If that happens we simply assume that the excess 
supply of the time $0$ good, denoted $\Omega_0 \geq 0$, 
is left to perish because no agents who want  those 
goods can also afford to buy them. Hence, the market
clearing conditions under the alternative 
equilibrium concept with a clearinghouse are
$$\EQNalign{ 
c_t^t + c_t^{t-1} &= y_t^t + y_t^{t-1}, 
\quad \forall t \geq 1;        \EQN olg_equib;a  \cr
c_0^0 + c_0^{-1} + \Omega_0 &= y_0^0 + y_0^{-1},
                                \EQN olg_equib;b \cr} $$
where $\Omega_0 \geq 0$.

Excess supply of time $0$
goods cannot occur in the chapter \use{recurge}
general equilibrium model. Consequently, it doesn't  matter
whether we compute a Walrasian equilibrium or instead use our
alternative equilibrium concept that entails a
clearinghouse for that economy. Both equilibrium 
concepts lead to the same unique equilibrium allocation 
that we characterized in chapter \use{recurge}.
``Walras' Law'' holds in that economy. Restated in
terms of our present setup, Walras' Law would
assert  that if all time $0$ markets for goods dated 
$t\geq 1$ clear, then  so must the remaining market $t=0$.
But we can and shall  show that  Walras' Law can be 
violated in an overlapping generations model.
So here it matters that we adopt our  alternative
equilibrium concept with its clearinghouse; this equilibrium concept lets us
 analyze a richer set of possible 
outcomes than encompassed within  a Walrasian equilibrium. 
Indeed, inspired by the additional 
outcomes that can  emerge under the alternative 
equilibrium concept with a clearinghouse, in later sections of this chapter we shall 
study sequential trading equilibria 
in which fiat money and unbacked government debt 
are   valued in equilibrium, outcomes that are not possible in the chapter \use{recurge} general 
equilibrium model.





%Lars%
%We use the definition of competitive equilibrium from 
%chapter \use{recurge}.
%Thus, we temporarily suspend disbelief and proceed in the 
%style of Debreu (1959) with time $0$ trading.
%\auth{Debreu, Gerard}%
%Specifically, we imagine that  at
%time $0$  an ``auctioneer'' or  ``clearinghouse''  posts 
%prices, monitors demands and supplies for goods in all 
%periods,
%and verifies that each person's single time $0$ budget 
%constraint is satisfied.
%An equilibrium price vector assures that for
%periods $t\geq 2$ aggregate endowments in that period are 
%consumed by living agents,  but it can turn out that 
%in period $1$ the aggregate endowment exceeds  aggregate 
%consumption so the time $0$ Arrow-Debreu  clearinghouse might
%end up with some of the time $1$ good not yet allocated.  
%Presumably, after trading concludes, the manager
% of the artificial clearing house could give  any such excess 
%supply of goods in period $1$  to the initial old
%generation without any effects on the equilibrium price 
%vector, since those old  agents optimally consume all
%their wealth in period $1$ and do not want to buy  goods
%in future periods. Reasons for our peculiar  treatment of
%period $1$  will emerge as we proceed.


\subsection{Agents' optimization}

Thus,
at time $0$, there are complete markets  in time $t \geq 0$ consumption goods with prices 
$\{q^0_t\}_{t\geq0}$.\NFootnote{As shown in 
chapter \use{recurge}, only relative prices are
determined in an equilibrium. For absolute prices, we
need to choose a numeraire. Here we have chosen
the good at time $0$ as the numeraire. Hence, using
our notation in chapter \use{recurge}, $q^0_t$ is the
price of a unit of the  good to be delivered at time $t$ in terms
of units of the good at time $0$.}
Agent $i$'s budget constraint is 
$$ \sum_{t=0}^\infty q^0_t c_t^i 
                     \leq \sum_{t=0}^\infty q^0_t y_t^i .
   \EQN olgbud $$
Letting  $\mu^i$ be a Lagrange multiplier attached to the budget constraint of consumer $i\geq 0$, the consumer's
first-order conditions are
$$\EQNalign{\mu^i q^0_i &= u'(c^i_i), \EQN olg2;a  \cr
            \mu^i q^0_{i+1} &= u'(c^i_{i+1}), \EQN olg2;b \cr
c_t^i & = 0  \ {\rm if} \ t \notin \{i, i+1\} .\EQN olg2;c \cr} $$
The consumer's Euler equation can be deduced from 
equations \Ep{olg2;a} and \Ep{olg2;b},
$$
{u'(c^i_{i+1}) \over u'(c^i_i)} = {q^0_{i+1} \over q^0_i} 
          \equiv \alpha_i = { 1 \over R_i },
                                          \EQN olg_Euler
$$
where $\alpha_i$ is the price of one good dated $i+1$ in 
terms of the good dated $i$.\NFootnote{If we had instead used 
our chapter \use{recurge}, $\alpha_t$ would 
have been written as $q^t_{t+1}$.}
Later, in a sequential 
trading equilibrium, $\alpha_t$ will be the inverse of the 
one-period gross interest rate $R_i$ between 
periods $i$ and $i+1$.

As for an  initial old agent $i=-1$  who is alive only
in the first period of the economy, the optimal consumption
choice is simply to consume his entire wealth in 
period $0$. Under our specialization of endowment
streams, the initial old consumer's budget constraint is
$q^0_0 c^{-1}_0 \leq q^0_0 y^{-1}_0$ and hence, the initial old agent's optimal
consumption choice is $c^{-1}_0 = y^{-1}_0$.





\subsection{Example equilibria}\label{sec:olg_stat}%

\noindent 
Let $\epsilon \in (0, .5)$.
Endowments are
$$ \eqalign{ y_i^i  & = 1 - \epsilon,  \ \forall i \geq 0, \cr
              y_{i+1}^i & = \epsilon, \ \forall i \geq -1, \cr
              y_t^i &=  0 \ {\rm otherwise.} \cr} 
                                            \EQN example1  $$
This economy has many equilibria. We describe two stationary
equilibria now and nonstationary equilibria later.  
An equilibrium is said to be a stationary equilibrium if
its allocation is stationary in the following sense.
\medskip
\noindent{\sc Definition:} An allocation is {\it stationary} if
$c^i_i = c_y$ and $c^i_{i+1} = c_o$ for all $i \geq 0$.
\medskip
\noindent Here the subscript $y$ denotes young and $o$ 
denotes old.
Note that we do not require that $c^{-1}_0 = c_o$.
Indeed, as explained above, equilibrium
consumption of the initial old agent $i=-1$ will
be $c^{-1}_0 = y^{-1}_0 = \epsilon$ in  both stationary and nonstationary
equilibria.
\index{equilibrium!multiple}

We can use a \idx{guess-and-verify method} to
confirm two stationary equilibria.
Specifically, we proceed by first guessing the 
equilibrium allocation for agents $i\geq 0$, 
which amounts to guessing $c_y$ and $c_o$ in a 
stationary equilibrium. Given that guess, we 
use Euler equation \Ep{olg_Euler} to deduce at 
which relative price $q^0_{i+1}/q^0_i$ such
a consumption profile $(c_y,c_o)$ is 
consistent with utility maximization by
agent $i\geq 0$. Next, given that deduced 
relative price and the guess of an 
allocation, we check that budget constraint 
\Ep{olgbud} holds with equality for each agent 
$i\geq 0$, which here becomes
$$ c_y + {q^0_{i+1} \over q^0_i} c_o 
=(1-\epsilon) + {q^0_{i+1} \over q^0_i} \epsilon,
 \EQN olgbud2 $$
where we have divided both sides of the 
budget constraint by $q^0_i$ and hence, 
changed the numeraire to be the good dated $i$.
Finally, we verify that markets clear
in accordance with equations \Ep{olg_equib}.

So by using this guess-and-verify method,
we can confirm the following two stationary
equilibria.


\medskip
\item{1.}  Equilibrium H: a high-interest-rate
equilibrium.  We guess a stationary equilibrium
allocation $c_y=c_o =.5$ and $c^{-1}_0 = \epsilon$. 
To be consistent with utility maximization, 
Euler equation \Ep{olg_Euler} must be satisfied
and hence, the relative price would have to be
$q^0_{i+1}/q^0_i = u'(.5)/u'(.5)=1=\alpha_H$. Given our 
normalization $q^0_0=1$, the sequence of absolute 
prices becomes $q_t^0 =1$ for all $t \geq 0$. 
We can then verify that budget constraint 
\Ep{olgbud2} holds with equality for each agent 
$i\geq 0$, and trivially for the initial old
agent $i=-1$ who simply consumes his endowment
in period $0$, $c^{-1}_0 =\epsilon$. 
Concerning market clearing conditions,
equation \Ep{olg_equib;a} holds since aggregate
consumption equals aggregate endowments 
in periods $t\geq 1$, while equation 
\Ep{olg_equib;b} yields 
excess supply of goods dated $0$, namely, 
$$ \Omega_0 = y^0_0 + y^{-1}_0 
    - c^0_0 - c^{-1}_0 
%%=(1-\epsilon) + \epsilon -.5 - \epsilon
= .5 - \epsilon > 0,
$$
where the quantity $\Omega_0>0$ of goods 
dated $0$ is left to perish since no agent
who wants those goods can afford to buy them.


\medskip
\item{2.}  Equilibrium L: a low-interest-rate
equilibrium. We guess a stationary equilibrium
allocation $c_y=1-\epsilon$, $c_o =\epsilon$,
and $c^{-1}_0 = \epsilon$. For that consumption
profile to be consistent with utility maximization,
the relative price would have to be
$q^0_{i+1}/q^0_i = u'(\epsilon)/u'(1-\epsilon)
= \alpha_L > 1$, 
i.e., Euler equation \Ep{olg_Euler} must hold
for all $i\geq 0$. Given our 
normalization $q^0_0=1$, the sequence of absolute 
prices becomes $q_t^0 = (\alpha_L)^t$ for all 
$t \geq 0$. Since all agents are simply 
consuming their endowments, 
$c^i_t = y_t^i$ for all $i, t$, budget constraints
\Ep{olgbud} and market clearing conditions
\Ep{olg_equib} are trivially satisfied, and 
there is no excess supply of goods dated $0$,
$\Omega_0=0$. This equilibrium is autarkic,
with equilibrium prices eradicating all incentives 
to trade.


\medskip
\noindent To understand why we call them high- and 
low-interest-rate equilibria, recall that the
relative price $\alpha_t$ in expression
\Ep{olg_Euler} would be the inverse of the
one-period gross interest rate $R_t$ in a
sequential trading equilibrium. Thus, the
interest rates associated with the high- 
and low-interest-rate equilibrium 
are $R_H=1$ and $R_L=(\alpha_L)^{-1}<1$,
respectively.
Soon we shall construct equilibria in which 
$R_t$ varies over time $t$.



%\medskip
%\item{1.}  Equilibrium H: a high-interest-rate
%equilibrium.  Set $q_t^0 =1 \ \forall t \geq 1$ and
%$c^i_i = c^i_{i+1} = .5$ for all $i \geq 1$ and $c^0_1 = 
%\epsilon$.
%   To verify
%that this is an equilibrium, notice that
%each consumer's first-order conditions are satisfied
%and that the allocation is feasible.  At the equilibrium 
%price vector
%$q_t^0$, extensive
%intergenerational trade  occurs
%at time $0$.   Constraint \Ep{olg3} holds with  equality
%for all $t \geq 2$ but with strict inequality for $t=1$; note %that some of the $t=1$ consumption good is left unconsumed.

%\medskip
%\item{2.}  Equilibrium L: a low-interest-rate
%equilibrium.  Set $q_1^0 = 1$, ${q^0_{t+1}  \over
%q^0_t} = {u'(\epsilon) \over u'(1 - \epsilon)} = \alpha > 1$.
%Set $c^i_t = y_t^i$ for all $i, t$.   This equilibrium is 
%autarkic,
%with equilibrium prices  eradicating all incentives to trade.

%\medskip
%\noindent To understand why we call them high- and low 
%interest rate equilibria,
%note that we can represent $q^0_t$ as
%$$ q_t^0 = q_1^0  \exp(-r_{1,2}) \cdots \exp(-r_{t-1,t}) $$
%where $R_{t,t+1}^{-1} = {\frac{1}{1+r_{t,t+1}}} \approx 
%\exp(- r_{t,t+1})$
%is the reciprocal of the gross interest rate between times 
%$t$ and $t+1$. 
%In Equilibrium H, $R_{t,t+1}^{-1} = 1$ for all $t \geq 1$ 
%while in equilibrium L ,
%$R_{t,t+1}^{-1} = \alpha > 1$ for all $t \geq 1$.
%Soon we shall construct  equilibria in which $R_{t,t+1}^{-1}$ 
%varies over time $t$.

\medskip
\subsection{Relation to welfare theorems}

Equilibrium H Pareto dominates
equilibrium L.  In equilibrium H, every generation after the initial
old one is better off and no generation is worse off than in equilibrium
L.
%As we shall explain more below,
%the equilibrium  1 Pareto dominates equilibrium 2.  Every generation
%after the initial old one is better off in the second equilibrium,
%and no generation is worse off.
The equilibrium H allocation
is peculiar because some of the time $0$ good is not consumed, leaving
 room to set up a giveaway program to the initial
old that makes them better off.  We shall soon learn how an institution of either fiat money or of perpetual government debt 
can accomplish this purpose.\NFootnote{See Karl Shell (1971) for an
investigation that characterizes why some competitive equilibria
in overlapping generations models fail to be Pareto optimal. 
As explained by Shell, ``the double infinity of traders and
dated commodities allows for competitive equilibria that are
not Pareto-optimal [and] the fact that generations do not 
meet is not essential.''}
%Shell also reviews earlier studies that had sought to
%understand how the welfare theorems seem to fail in the
% overlapping generations structure.}
%NOT EXACTLY, SHELL'S 2ND HALF FOCUSES ON SOMETHING ELSE
\auth{Shell, Karl}

With equality between supply and demand in all
periods, equilibrium L is also a Walrasian equilibrium.
But, evidently, this competitive equilibrium
fails to satisfy
  the assumptions needed to deliver the
   first fundamental theorem of welfare economics that isolates conditions under which
 a competitive equilibrium allocation is Pareto
optimal.\NFootnote{See Mas-Colell, Whinston, and Green (1995) and Debreu (1954).
\auth{Mas-Colell, Andreu} \auth{Whinston, Michael D.} \auth{Green, Jerry R.}%
\auth{Debreu, Gerard}
}
  The condition that is violated by  equilibrium L
is the assumption that the value of the aggregate endowment
at the equilibrium prices is finite.\NFootnote{Note
that if the horizon of the economy were finite, then the counterpart of
equilibrium H would not exist and the  allocation of the
counterpart of equilibrium L would be Pareto optimal.}

\subsection{Nonstationary equilibria}\label{sec:olg_nonstat}%

\noindent 
Our example economy has more  equilibria.
To construct  them, we use an offer curve to summarize preferences and optimal consumption
decisions.   We describe a graphical apparatus proposed
by David Gale (1973) and used  to great advantage by
William Brock (1990).
\auth{Gale, David}%
\auth{Brock, William A.}%

\medskip
\noindent{\sc Definition:}  Agent $i$'s {\it offer curve} is
a locus of $(c_i^i, c^i_{i+1})$ that solves
$$ \max_{\{c_i^i, c^i_{i+1}\}} U^i({\bf c}^i)  $$
subject to
$$ c_i^i + \alpha_i c_{i+1}^i \leq y_i^i + \alpha_i y^i_{i+1} .$$
Here $\alpha_i \equiv {q^0_{i+1}\over q^0_i } $,
 the reciprocal
of the one-period gross rate of return from period $i$ to $i+1$,  is 
a parameter that we shall use to sweep out a $(c_i^i, c_{i+1}^i)$ locus.
\medskip
Evidently, points $(c_i^i, c_{i+1}^i)$ on
 the offer curve simultaneously solve the following pair of equations:
$$\EQNalign{c^i_i + \alpha_i c^i_{i+1} 
        & = y_i^i + \alpha_i y^i_{i+1} \EQN offer1;a  \cr
            {u'(c_{i+1}^i) \over u'(c^i_i)} & = \alpha_i   
                                        \EQN offer1;b  \cr}$$
for  different values of the relative price $\alpha_i > 0$.  
%for  indifference curve slope $\alpha_i > 0$.  
We denote the offer curve  swept out by $\alpha_i > 0 $  by
$$ \psi(c^i_i, c^i_{i+1}) = 0 .$$
%Evidently, an offer curve lies inside the indifference curve that goes through the person's endowment point.


The offer curve is illustrated in Figure \Fg{graph12f}. %Figure 8.1.
We trace it out by varying $\alpha_i$ in the consumer's problem and
reading off tangency points between the consumer's indifference
curve and the budget line.  
So at each point on the offer curve, there
is an indifference curve whose tangent at
that point  is a straight line that can
be extended to go through the consumer's 
endowment point.
Hence, the locus of the offer curve depends
on the endowment vector and  lies above the 
indifference curve through
the endowment point.  
Since the consumer can always choose to 
consume his endowment, an  optimally
chosen consumption $(c_i^i, c^i_{i+1})$  could 
never lie on an indifference curve lower than 
the indifference curve through the endowment
point. 


%Lars%
%The resulting locus depends
%on the endowment vector and  lies above the indifference 
%curve through
%the endowment vector.  At the intersection between the offer 
%curve and a straight
%line through the endowment point, the straight line is 
%tangent to an indifference curve.


By relabeling axes, in the same graph we can plot a consumption possibility frontier
$$ c_i^i + c_i^{i-1} \leq y_i^i + y_i^{i-1}    \EQN Cfront$$
that bounds feasible allocations at time $i$, which we
call the {\it feasibility line}.  Note that while the offer
curve restricts a life-time consumption
profile $(c_i^i, c_{i+1}^i)$ for one person born at time $i$, the feasibility line involves 
pairs $(c_i^i, c_i^{i-1})$ of consumptions assigned to a young and old person at time $i$.\NFootnote{Given our assumptions 
on preferences and endowments,
the conscientious reader will note that  Figure \Fg{graph12f} %Figure 8.1
appears  distorted because  the offer curve really ought to intersect the feasibility line
along the 45 degree line  with $c^t_t=c^t_{t+1}$, i.e., at the allocation affiliated with equilibrium H above.}


%It depends implicitly on the endowment. We trace it out by
%varying $\alpha$ in the consumer's problem and reading
%tangency points between the consumer's indifference curve
%and the budget line. The offer curve lies
%above the indifference curve through the initial endowment (because
%trade is voluntary).  And once again, by construction,
%an indifference curve through the
%intersection of the offer curve and
%a straight line between   a point on the offer curve and the endowment
%must be tangent to that line (this follows from \Ep{offer1}).
%  See figure 1.

%%%%%%
%$$
%\grafone{graph12.eps,height=2.5in}{{\bf Figure 8.1} The offer curve and
%feasibility line.}
%$$
%%%%%%%

\midfigure{graph12f}
\centerline{\epsfxsize=3truein\epsffile{graph12.eps}}
\caption{The offer curve and feasibility line.}
\infiglist{graf12f}
\endfigure

\medskip \auth{Gale, David}%
Following Gale (1973), for the important special case in which $(y_i^i, y_i^{i-1}) = (y_y, y_o)$ for all $i \geq 0$,
 we can use the offer curve and
a straight line depicting feasibility in the $(c^i_i, c^{i-1}_i)$ plane
to  compute equilibrium allocations and prices.  The assumption 
of stationary endowments
%that  $(y_i^i, y_i^{i-1}) = (y_y, y_o)$ for all $i \geq 0$
means that the feasibility line is the same in
all periods $i \geq 0$. As explained earlier, the initial 
old agent consumes his endowment in an equilibrium,
$c^{-1}_0 = y_o$. As for the continuation of the equilibrium 
allocation, there is a multiplicity of possible
outcomes since $y_o<y_y$. Specifically, we can initiate an
equilibrium computation by choosing
$c^0_0 \in ((y_y+y_o)/2,y_y)$.\NFootnote{By setting 
$c^0_0$ equal to one of the two boundary points, 
$(y_y+y_o)/2$ or $y_y$, we could recover one of the two section \use{sec:olg_stat} stationary equilibria.}
As we shall see, for each  different $c_0^0$, a different  competitive equilibrium prevails.
Thus, we  can  find an equilibrium allocation by solving  a pair of nonlinear difference equations.
  For $i \geq 0$, the equations
are
%\NFootnote{By imposing equation \Ep{equilibr1;b} with
%equality, we are possibly including a giveaway program
%to the initial old.}
$$\EQNalign{  \psi(c^i_i, c^i_{i+1}) & = 0, \EQN equilib1;a \cr
              c^{i+1}_{i+1} + c^i_{i+1} & = y_y + y_o. \EQN equilibr1;b \cr}$$
Equation \Ep{equilib1;a} states that the allocation to an agent
born at time $i$ is on the offer curve. 
Equation \Ep{equilibr1;b}
asserts that the allocation at time $i+1$  is 
feasible.\NFootnote{Feasibility in 
period $i=0$ is assured by our choice of $c^0_0<y_y$ and
the initial old agent consuming his endowment in any
equilibrium, $c^{-1}_0=y_o$, i.e., equation \Ep{olg_equib;b}
is satisfied with $\Omega_0 > 0$.\label{note:feasible0}}
So we take $c_0^0$ as an initial condition and then compute an equilibrium allocation via the following steps:


\medskip
\item{({\bf a})} Solve agent $0$'s offer curve  \Ep{equilib1;a} for $c^0_1$.

\medskip
\item{({\bf b})} Solve feasibility condition \Ep{equilibr1;b} for $c_1^1$.

\medskip
\item {({\bf c})} Iterate on these steps indefinitely and hope for convergence.


%\medskip
%\item{({\bf d})} An optional last step: solve feasibility 
%condition \Ep{equilibr1;b} perhaps for a set of  $c_1^0$'s that 
%the clearing house 
%could give to the initial old if it wishes.


\medskip

\noindent After an allocation has been computed, a time $0$ 
trading equilibrium price system can be computed from
$$ q^0_{i+1} = q^0_i {u'(c^i_{i+1}) \over u'(c^i_i) } $$
for all $i \geq 0$, given the numeraire $q^0_0 =1$.






\subsection{Computing equilibria}\label{sec:olg_time0log}%

We now describe David Gale's (1973) equilibrium computation machine that consists of a graph like that
%  Figure 8.2 is David Gale's (1973) machine for
%computing an equilibrium.
%It works as follows.
 illustrated
in Figure \Fg{graph22f}.
% %Figure 8.2,
%which reproduces a version of a graph
%of David Gale (1973).
Notice that  $y_t^t$ and  $c_t^t$  are on the ordinate axis while
$y_{t+1}^t$ and both $c_{t+1}^t$ and $c_{t}^{t-1}$ are on the horizontal axis.
The straight  feasibility line with slope $-1$ goes through endowment point $(y_t^t, y^t_{t+1})$ and plots feasible
pairs   $(c_t^t, c_t^{t-1})$ that can be allocated to people alive at time $t$.  The curved offer curve also goes through the endowment point $(y_t^t, y^t_{t+1})$
and plots pairs $(c_t^t, c^t_{t+1})$ that a person born at $t$ optimally chooses as a function of the relative price $\alpha_t$ implicitly
sweeping out the offer curve.   

%Lars%\medskip
%\noindent{\it Example 1:}  \quad  Plain vanilla model with 
%stationary endowment pattern:
%  This is the situation illustrated
%in Figure \Fg{graph22f}.  

Applying the section \use{sec:olg_nonstat} algorithm, we start with
a proposed $c_0^0$, a time $0$ allocation to the initial young.
%Lars%Then use the feasibility line to  find the {\it maximal\/}
%feasible value for $c_0^1$, the time $1$  allocation to    the
%initial old.  In the Arrow-Debreu equilibrium, the allocation to
%the initial old will be less than this maximal value, so that 
%some of the time $1$ good is thrown away.  This follows because
% the budget constraint of the initial old, 
%$q_1^0( c^0_1 - y^0_1) \leq 0$, implies
%that $c^0_1 = y^0_1$.\NFootnote{Soon we shall discuss another 
%market structure that
%avoids throwing away any of the initial endowment by augmenting 
%the endowment of the
%initial old with a particular zero-dividend infinitely durable 
%asset.}
The candidate time $0$ equilibrium allocation is feasible as explained in  footnote \use{note:feasible0}, but  the
time $0$  young will choose $c_0^0$ only if the price $\alpha_0$ is such that
$(c^0_0, c^0_1)$ lies
on the offer curve.      Therefore, we choose $c^0_1$ from
the point on the offer curve that cuts a vertical
line through  $c^0_0$.   Then we proceed to find $c^1_1$ from
the intersection of a horizontal line through $c^0_1$
and the  feasibility line.   We continue recursively in this way,
choosing $c_i^i$ at the intersection of the feasibility line
with a horizontal line through $c^{i-1}_i$, then
choosing $c_{i+1}^i$ at the intersection of a vertical
line through $c_i^i$   and the offer curve.    We can
construct a sequence of $\alpha_i$'s from the slope
of a straight line through the endowment point
and the sequence of $(c_i^i, c^i_{i+1})$  pairs that
lie on the offer curve.

 If the offer curve has the shape drawn in Figure \Fg{graph22f}, any %Figure 8.2, any
$c_0^0$ between the 
%upper and lower intersections
two intersections
of the offer curve and the feasibility line is
an equilibrium setting of $c_0^0$. Each such $c_0^0$ is
associated with a distinct allocation 
$\{c^i_i, c^i_{i+1}\}_{i\geq 0}$ and
price system $\{\alpha_i\}_{i\geq 0}$, while  
consumption of the initial old agent is the
same across all equilibria, $c^{-1}_0 = y_o$.
If $c^0_0$ is chosen to be the lower boundary
point for that set, the economy will simply 
remain in that {\it high\/}-interest-rate stationary 
equilibrium. For all other choices of $c^0_0$,
the economy converges to the {\it low\/}-interest-rate
stationary equilibrium allocation and interest rate.

%Lars%
%all but
%one of them converging to the {\it low\/}-interest-rate
% stationary equilibrium allocation and interest rate.


%%%%%%%%
%$$
%\grafone{graph22.eps,height=2.5in}{{\bf Figure 8.2} A nonstationary
%equilibrium allocation.}
%$$
%%%%%%

\midfigure{graph22f}
\centerline{\epsfxsize=3truein\epsffile{graph22.eps}}
\caption{A nonstationary equilibrium allocation.}
\infiglist{graph22f}
\endfigure

\medskip
\noindent{\it Log utility:} \quad 
Suppose that $u(c) = \ln c$ and that the
endowment is described by equations \Ep{example1}
with $\epsilon \in (0, .5)$.   
After solving equations \Ep{offer1} for that 
specification, the offer curve is given by the 
recursions
$$\EQNalign{
   c^i_i & = .5 (1 - \epsilon + \alpha_i \epsilon) ,
                                        \EQN offerlog;a  \cr
   c^i_{i+1} & = {1 \over \alpha_i} c_i^i .  \EQN offerlog;b  \cr}
$$
Hence, feasibility at $i\geq 1$ and the offer curves then
imply
$$ {1 \over 2} (1 -\epsilon + \alpha_i \epsilon)  
  +  {1 \over 2 \alpha_{i-1}} 
     (1 - \epsilon + \alpha_{i-1} \epsilon) =1. \EQN alphadiff 
$$
This implies the difference equation
$$ \alpha_i = \epsilon^{-1} - {\epsilon^{-1} -1 \over \alpha_{i-1}}.
  \EQN alphadiff1 $$
To find the stationary equilibria of this economy, we
let $\alpha_i = \alpha$ for all $i$ in 
difference equation \Ep{alphadiff1} and after
rearranging,
$$ \alpha^2 - \epsilon^{-1} \alpha 
   + \left( \epsilon^{-1} -1 \right) = 0.
                                           \EQN alphadiffstat 
$$
Using the quadratic formula, the equilibrium prices in
two stationary equilibria are:
$$
\alpha = { \epsilon^{-1}  \pm \sqrt{ \epsilon^{-2} 
           - 4 \left( \epsilon^{-1} -1 \right) } \over 2 }
       =  { \epsilon^{-1} \pm \sqrt{ 
             \left( \epsilon^{-1} -2 \right)^2 } \over 2 }
   = \cases{  1 ,  \cr
%   \noalign{\vskip.3cm}  \cr
   {\displaystyle  1-\epsilon \over \displaystyle \epsilon} \cr}
                                    \EQN alphadiffstatsol
$$
The two equilibria are the log-utility versions
of the two stationary equilibria in
section \use{sec:olg_stat}, 
$\alpha_H  %=u'(.5)/u'(.5)
=1$ and
$\alpha_L=u'(\epsilon)/u'(1-\epsilon)=(1-\epsilon)/\epsilon >1$. 
To initiate the computation of a nonstationary equilibrium
using difference equation \Ep{alphadiff1}, we choose 
$\alpha_0 \in (\alpha_H, \alpha_L)$, 
which is tantamount to choosing 
$c^0_0 \in %((y_y+y_o)/2,y_y) = 
(.5, 1-\epsilon)$ when
using the section \use{sec:olg_nonstat} algorithm. 
The limiting value of $\{\alpha_i\}$ for any nonstationary
equilibrium is $\alpha_L$. Hence, the autarkic equilibrium
$\alpha_L$ is stable, while the stationary equilibrium
$\alpha_H = 1$ is not.



%Lars%
%Then the offer curve
%is given by the recursions $c^i_i = .5(1 - \epsilon + \alpha_i
%\epsilon), c^i_{i+1} = \alpha_i^{-1} c_i^i$.
%Let $\alpha_i$ be the gross rate of return facing   the young
%at $i$.   Feasibility at $i$ and the offer curves then
%imply
%$$ {1 \over 2 \alpha_{i-1}} (1 - \epsilon + \alpha_{i-1} 
%\epsilon)
%    + .5 (1 -\epsilon + \alpha_i \epsilon) =1. \EQN alphadiff $$
%This implies the difference equation
%$$ \alpha_i = \epsilon^{-1} - {\epsilon^{-1} -1 \over 
%\alpha_{i-1}}.  \EQN alphadiff1 $$
%See Figure  \Fg{graph22f}. %Figure 8.2.
% An equilibrium $\alpha_i$ sequence must
%satisfy  equation \Ep{alphadiff} and have $\alpha_i > 0$ for
%all $i$.   Evidently, $\alpha_i =1$ for all $i\geq 1$ is
%an equilibrium $\alpha$ sequence.   So is any $\alpha_i$ 
%sequence that  satisfies
%equation \Ep{alphadiff} and $\alpha_1 \geq 1$; $\alpha_1 < 1$ 
%will not work
%because equation
%\Ep{alphadiff} implies that the tail of $\{\alpha_i\}$ is an 
%unbounded
%negative sequence.  The limiting value of $\alpha_i$ for any
%$\alpha_1 >1  $ is ${1-\epsilon \over \epsilon} = u'(\epsilon) /
%u'(1-\epsilon)$,
%which is the interest factor associated with the stationary
%autarkic equilibrium.    Figure \Fg{graph22f} suggests that   
%the stationary
%$\alpha_i=1$ equilibrium is not stable, while the autarkic 
%equilibrium is.





\section{Sequential trading}
  We now alter the trading arrangement to bring it
into line with more common presentations of the overlapping
generations model.   We abandon the time $0$, complete markets
 trading  arrangement and replace it with sequential
trading. In our model with two-period lived agents, the 
only equilibrium is then autarky since there is no double
coincidence of wants. While a young agent might want to
buy or sell goods in exchange for 
goods next period, a current old agent will not be
around next period.
Hence, the unique equilibrium price system  eradicates all incentives to trade,
$\alpha_t= u'(y^t_{t+1})/u'(y^t_t)$ for all $t\geq 0$.

To create opportunities for trade across
generations, we introduce a government
that at time $0$ issues a zero-dividend infinitely
durable asset in the form of fiat money or unbacked
government debt. Under the auxiliary assumption 
that the government gives this asset 
to the initial old agent, we will show that all the 
equilibria of the time $0$ trading economy
can reemerge as sequential trading equilibria, but
with a  difference. Any goods that the clearinghouse 
had earlier left to perish at time $0$, 
$\Omega_0\geq 0$, will no longer be wasted but be
consumed by the initial old agent.
We study sequential trading
equilibria under the common
notion of a Walrasian auctioneer who announces
equilibrium prices at which the goods and asset
markets clear in every period.

The assumption that the government-issued asset 
is given to the initial old generation 
is critical for delivering  equivalence of 
outcomes under time $0$ trading and sequential 
trading with respect to equilibrium prices and
the allocation $\{c^t_t, c^t_{t+1}\}_{t\geq0}$
(while leaving aside the value of $c^{-1}_0$).
Any giveaway program to the 
initial old generation has no effects on the 
equilibrium price system, since 
those old agents optimally consume all their
wealth in period $0$. If the assets had
instead been given to the initial young generation,
it would have changed their offer curves and hence, 
affected equilibrium prices.


%Lars%
% We abandon the time $0$, complete markets
% trading  arrangement and replace it with sequential
%trading in which a durable asset, for example,  either 
%government debt or unbacked fiat money, 
%is passed from old to young.
%Some cross-generation transfers occur with voluntary exchanges,
% while others are engineered by government
%tax and transfer programs.


\subsection{Government budget constraint}
In the next sections, we will activate different items
in the following government budget constraints: 
$$ {M_t - M_{t-1} \over  p_t} + \tau^{t-1}_t + \tau^t_t 
+ {B_{t+1} \over R_t} = g_t + B_t,  \quad \forall t \geq 0, \EQN Gov_bc $$
where $M_t$ is the nominal quantity of money at the end of period $t$,
$p_t$ is the price level in period $t$, $\tau^i_t$ is a lump-sum
tax levied in period $t$ on the representative agent born in 
period $i\in\{t-1, t\}$ (a negative tax is a lump-sum transfer), 
$B_{t+1}$ is
one-period real bonds issued by the government in period $t$ which
pay $B_{t+1}$ units of goods upon maturity in period $t+1$, $R_t$
is the gross real one-period return on those bonds, and
$g_t$ is government purchases of goods at time $t$. Thus,
the left side of budget constraint \Ep{Gov_bc} sums up
government funds -- real seigniorage from money creation,
tax revenues, and the issue of government bonds -- which are
used to pay for goods purchases and the maturing debt on
the right side of \Ep{Gov_bc}. The economy starts
at time $t=0$ and the initial conditions are 
$M_{-1}= B_0 = 0$. We assume that government purchases,
$\{g_t\}_{t \geq 0}$, do not give utility to agents.



\section{Fiat money}\label{sec:Samuelson101}%
\auth{Samuelson, Paul}%
 In  Samuelson's (1958)  model, trading occurs
 sequentially through a medium of exchange, an
inconvertible (or ``fiat'') currency.
   Preferences and  endowments
are as described above, with one important additional component of the
endowment.
At date $t=0$, old agents are endowed in the aggregate
with $M >0$ units of intrinsically worthless currency that
no one has promised to redeem  for goods.
Thus,  currency is  not ``backed'' by another asset.  But 
Samuelson showed that there  exists a   system of expectations that  can give positive value to
an unbacked currency. Currency will be valued today if people expect it
to be valued tomorrow.  Thus, Samuelson studied a situation in which
currency is backed
by expectations without promises.
\index{expectations!without promises}

The described monetary economy of Samuelson (1958)
translates into a single government budget constraint
per our formulation \Ep{Gov_bc}:
$$ {M_0 \over  p_0} + \tau^{-1}_0 = 0,  \EQN Gov_bc_M $$
where $M_0=M$. That is, the government hands over the
currency $M$ to the initial old agents in period $0$. The
real value of that lump-sum transfer depends on the 
equilibrium price level $p_0> 0$, that is, 
$-\tau^{-1}_0 = M / p_0 \geq 0$. 
Thereafter, the government takes no further action
in the economy.

 

At each date $t \geq 0$, a young agent purchases
$m^t_t$ units of currency at a price of $1/p_t$ units
of the time $t$ consumption good.  Here $p_t> 0$ is
the time $t$ price level. At each $t\geq 0$, each
old agent exchanges his holdings of currency for the time $t$
 consumption good.    The budget constraints of a young
agent born in period $t \geq 0$ are
$$ \EQNalign{ c_t^t + {m_t^t \over p_t} & \leq y_t^t, \EQN samuels1 \cr
            c_{t+1}^t & \leq y^t_{t+1} + {m_t^t   \over p_{t+1} },
                \EQN samuels2 \cr
           m_t^t & \geq 0. \EQN samuels3 \cr} $$
%If $m_i^i \geq 0$
If money is valued and, hence, held by agents, inequalities \Ep{samuels1} and \Ep{samuels2} imply
$$ c_t^t + c^t_{t+1} \left( {p_{t+1} \over p_t}\right) \leq y_t^t + y^t_{t+1}
          \left(  {p_{t+1} \over p_t}\right) . \EQN samuels4 $$
Provided that we set
$$ {p_{t+1} \over p_t} = \alpha_t = {q^0_{t+1} \over q^0_t },$$
 this budget set is identical with the one implied by equation \Ep{olgbud}
for agent $i=t\geq0$. However, the budget set for an initial old
agent, $i=-1$, differs from that implied by equation \Ep{olgbud} and
is now given by  
$$ c_0^{-1} \leq y_0^{-1} -\tau^{-1}_0 . \EQN samuels4_ini $$
where $-\tau^{-1}_0 = M / p_0 > 0$ by the government's
budget constraint \Ep{Gov_bc_M}.


  We use  the following definitions:
\medskip
\noindent{\sc Definition:}  A nominal price  sequence is a strictly positive
sequence $\{p_t\}_{t \geq 0}$.
\medskip
\noindent{\sc Definition:}  An equilibrium with valued fiat money
is a feasible allocation and a nominal price sequence with $p_t < +\infty$
for all $t\geq 0$
 such
that given the price sequence, the allocation solves the agent's
problem for each generation $i\geq -1$.
\medskip


\noindent The qualification that $p_t < + \infty$ for all $t$  means that
fiat money is valued. We call an equilibrium with valued fiat money
a `monetary equilibrium'.  If $p_t = + \infty$, we   call it a `nonmonetary equilibrium'.



\subsection{Computing equilibria with valued fiat currency}\label{sec:Meqlog}%
   Summarize a  consumer's optimal decisions with
a saving function
$$ y_i^i - c_i^i = s(\alpha_i; y_i^i, y_{i+1}^i). \EQN savfn $$
Then the equilibrium conditions  are
$$ \EQNalign{ {M  \over p_i} & = s(\alpha_i; y_i^i, y_{i+1}^i), \EQN sameq1;a \cr
                \alpha_i & = {p_{i+1} \over p_i}, \EQN sameq1;b \cr }$$
where it is understood that
$ c_{i+1}^i = y_{i+1}^i + {M \over p_{i+1}} $.
Equation \Ep{sameq1;a} states that at time $i$   net  saving of generation $i$ (the expression on the right side) equals
 net dissaving of generation $i-1$ (the expression on the left side).  To compute an equilibrium, we solve the difference equations
\Ep{sameq1} for $\{p_i\}_{i\geq 0}$,  then get the allocation
from the consumer's budget constraints evaluated at equality
at an equilibrium level of real balances $M/p_i$.

As an example, suppose that $u(c) = \ln(c)$,
and that $(y^i_i, y^i_{i+1}) = (w_1, w_2)$ with $w_1 > w_2$.
  The saving function is
$ s(\alpha_i; w_1, w_2) = .5(w_1 - \alpha_i w_2)$.   Then
equation \Ep{sameq1;a} becomes
$$ .5(w_1 - w_2 {p_{t+1} \over p_t}) = {M \over p_t} $$
or
$$ p_t = 2 M/ w_1 + \left({w_2 \over    w_1}\right) p_{t+1} .\EQN samdiff1 $$
This is a difference equation whose solutions
with positive price levels are
$$ p_t = {2 M \over w_1 (1 - {w_2 \over w_1})} + \kappa
  \left({w_1\over w_2}\right)^t, \EQN samsol1 $$
for any scalar $\kappa \geq 0$.\NFootnote{See  appendix \use{appa1}
from chapter \use{timeseries}.}
The solution  $\kappa=0$ is the unique stationary solution.  Solutions
with $\kappa >0$ have uniformly higher price levels than the $\kappa=0$ solution,
and have  values of currency approaching zero in the limit as $t \rightarrow +\infty$.

It is instructive to compare the multiple monetary equilibria implied by 
equation \Ep{samsol1}, indexed by $\kappa \in [0,\infty)$, to the time 0 
trading equilibria computed in section \use{sec:olg_time0log}. The 
stationary monetary equilibrium with a constant price level, $\kappa =0$, 
is the counterpart to the high-interest-rate stationary equilibrium under 
time 0 trading, $\alpha_H$. 
Since the two stationary equilibria share the same price of one good dated
$t+1$ in terms of goods dated $t$, $p_{t+1} / p_t = 1 = \alpha_H$, the 
consumption choices of agents born in period $t\geq 0$ will also be the same.
In particular, given $\kappa =0$, the constant real money supply implied by 
equation \Ep{samsol1} is $M/p_t = (w_1-w_2)/2$ for all $t\geq 0$ and, hence, 
the consumption allocation in the stationary monetary equilibrium of the 
representative agent born in period $t\geq 0$ is
$c_t^t = c_{t+1}^t = (w_1 + w_2)/2$.
%$$ \EQNalign{ c_t^t &= w_1 - {M  \over p_t} = {w_1 + w_2 \over 2},\EQN samelnc;a
%                          \cr
%          c_{t+1}^t &= w_2 + {M  \over p_{t+1}} = {w_1 + w_2 \over 2}. \EQN samelnc;b \cr}$$
This perfect consumption smoothing corresponds to $c_t^t=c_{t+1}^t=.5$ for $t\geq 0$ 
in the high-interest-rate stationary equilibrium under time 0 trading, given
our earlier parameterization $y_t^t=1-\epsilon$ and $y_{t+1}^t=\epsilon\in (0, .5)$. 
(Though, note that the monetary equilibrium and the time $0$ trading equilibrium
differ with respect to the consumption of the initial old agents in period $0$, as 
we discussed above and will revisit below.)

Another similarity of the stationary monetary equilibrium and the 
high-interest-rate stationary equilibrium under time 0 trading is that these
equilibria are not stable. In fact, all other equilibria converge to autarky,
i.e., the low-interest-rate stationary equilibrium under time 0 trading, $\alpha_L$.  
To confirm that the relative price $p_{t+1} / p_t$ in all other monetary
equilibria, $\kappa\in(0,\infty)$, converges to 
$\alpha_L=y_t^t/y_{t+1}^t$ as derived in section \use{sec:olg_time0log}, 
%%(1-\epsilon)/\epsilon$ in section \use{sec:olg_time0log},
we compute
$$
\lim_{t\to\infty} {p_{t+1} \over p_t} \Big{\vert}_{\kappa\in(0,\infty)}
= \lim_{t\to\infty} { { {2 M \over w_1 (1 - {w_2 \over w_1})} + \kappa
  \left({w_1\over w_2}\right)^{t+1} } \over { {2 M \over w_1 (1 - {w_2 \over w_1})} + \kappa
  \left({w_1\over w_2}\right)^t } } \Bigg{\vert}_{\kappa\in(0,\infty)}
%= \lim_{t\to\infty} { {  \kappa
%  \left({w_1\over w_2}\right)^{t+1} } \over {  \kappa
%  \left({w_1\over w_2}\right)^t } } \Bigg{\vert}_{\kappa\in(0,\infty)}
= {w_1\over w_2}.
$$




\subsection{Equivalence of equilibria}
%   We briefly look back at  equilibria with
%time $0$ trading and note that   equilibrium allocations are
%the same under time $0$ and sequential trading.
As illustrated by the example above, there is a particular equivalence 
of time 0 trading equilibria with a clearinghouse and sequential
trading with fiat money.  
Thus, the following proposition
asserts that  with an adjustment to the endowment and the consumption allocated
to the initial old, a competitive equilibrium
allocation with time $0$ trading
is a competitive equilibrium  allocation in the fiat money economy with
sequential trading.
\medskip
\noindent{\sc Proposition:}  Let $\overline {\bf c}^i$ denote a competitive equilibrium
 allocation with time $0$ trading
  and suppose that it satisfies
$\overline c^0_0 < y^0_0$.
There exists
an equilibrium allocation  of the monetary economy with sequential trading 
 that satisfies $c^i_i = \overline c^i_i$, $c^i_{i+1} = \overline
c^i_{i+1}$ for $i \geq 0$.

\medskip
\noindent{\sc Proof:}  Take the competitive equilibrium allocation
and  price system under time 0 trading 
and form $ \alpha_i = q^0_{i+1}/q^0_i$.  Set
$m_i^i/p_i = y_i^i - \overline c_i^i$. Set $m_i^i = M$ for all $i \geq 0$,
and   determine $p_0$ from ${M \over p_0} = y_0^0  - \overline c_0^0$.
This last equation determines a positive initial price level
$p_0$ provided that $y_0^0 - \overline c_0^0 >0$.  Determine
subsequent price levels from $p_{i+1} = \alpha_i p_i$.
Assign an allocation to the  initial old: % according to
$c^{-1}_0 = y^{-1}_0 + {M \over p_0} = y^{-1}_0 +(y_0^0 - \overline c_0^0)$.  \qed

\medskip
In a monetary equilibrium, time $t$ real balances
equal both the per capita {\it saving\/} of the young and the
per capita {\it dissaving\/} of the old. To be a monetary equilibrium, both quantities must
be positive for all $t \geq 0$.


A converse of the proposition is true.
\medskip
\noindent{\sc Proposition:}  Let    $\overline {\bf c}^i$ be an
equilibrium allocation for the fiat money economy.    Then
there is a competitive equilibrium with time $0$ trading
with the same allocation, provided that the endowment
of the initial old is augmented with an appropriate  transfer
from the clearinghouse.
\medskip

To  verify this proposition, we have to construct
the required transfer from the clearinghouse to the initial
old.  Evidently, it is $y^0_0 - \overline c^0_0$.  We invite
the reader to complete  the proof.



%To  verify this proposition, we have to construct
%the required transfer from the government to the initial
%old.  Evidently, it is $B =  y^1_1 - \overline c^1_1$.   The
%government awards  claims of this amount to the initial
%old.  These claims are denominated in units of time $1$ consumption
%goods, and are never to be redeemed.   By the mere act of issuing
%these claims,
%the government  creates wealth.      The government
%never has to  tax   to give these claims value.
%The market simply assigns positive  value  to them.
%This seems mysterious, because the government bonds
%are backed by no `fundamentals'.
%  That government bonds are valued without backing is
%closely connected to the fact that there is an
%equilibrium of the economy without government debt
%in which the equilibrium interest rate is below the
%growth rate of population.
%For now, we simply invite the
%reader to verify that the proposition is true.
%Later we'll shed more light on how such unbacked claims
%acquire value.


 \section{Unbacked government debt}
We will now map the fiat monetary equilibria of the preceding section 
into equilibria with unbacked government debt. Instead of emitting  
fiat currency in period $0$, the government finances the lump-sum transfer 
$-\tau^{-1}_0 >0$ to the initial old agents in period $0$ by 
issuing real (indexed) bonds. Specifically, let $B_{t}$ be the quantity of 
one-period bonds issued by the government in period $t-1$, 
which promise to pay $B_{t}$ units of goods upon maturing in period $t$. 
The price of those bonds in period $t-1$ is the inverse of the gross 
real one-period interest rate $R_{t-1}$. In each period $t \geq 1$, 
the government redeems the maturing debt by issuing new one-period bonds. 
Thus, the debt is unbacked since there are never any taxes levied to service 
the debt. The government budget constraints are 
$$\EQNalign{\tau^{-1}_0 + {B_1 \over R_0}&= 0,    \EQN Gov_bc_B;a \cr
                     {B_{t+1} \over R_t} &= B_t,  \quad \forall t \geq 1. 
                                                  \EQN Gov_bc_B;b \cr} $$
We solve equations \Ep{Gov_bc_B} forward to some date $T\geq 1$:
$$ -\tau^{-1}_0 = {B_1 \over R_0} = {1 \over R_0}\, {B_2 \over R_1}
= {1 \over R_0 R_1} \,{B_3 \over R_2 } = \ldots 
= {1 \over \prod_{j=0}^{T-1} R_j}\, {B_{T+1} \over R_{T} }   \EQN Gov_bc_iter
$$
and, hence, government indebtedness in period $T\geq 1$ equals
$$ {B_{T+1} \over R_{T} } = -\tau^{-1}_0 \, \prod_{j=0}^{T-1} R_j > 0, \EQN Gov_debt
$$
which arises from the initial borrowing at time 0, $B_1/R_0=-\tau^{-1}_0$,
and compounded interest.

At each date $t \geq 0$, a young agent purchases $b^t_{t+1}$ 
units of newly issued one-period government bonds at a price of 
$1/R_t$ units of the time $t$ consumption good. When the bonds
mature next period $t+1$, the government pays $b^t_{t+1}$ units of the 
time $t+1$ consumption good to the agent as a bond holder. 
The budget constraints of a young agent born in period $t \geq 0$ are
$$ \EQNalign{ c_t^t + {b_{t+1}^t \over R_t} & \leq y_t^t, \EQN bc_B;a \cr
            c_{t+1}^t & \leq  y^t_{t+1} + b_{t+1}^t.      \EQN bc_B;b \cr} $$
Inequalities \Ep{bc_B;a} and \Ep{bc_B;b} imply
$$ c_t^t + {c^t_{t+1} \over R_t} \leq y_t^t + {y^t_{t+1} \over R_t}. 
                                                           \EQN bc_B_pv $$
The budget constraint of an initial old agent, $i=-1$, is
$$ c_0^{-1} \leq y_0^{-1} -\tau^{-1}_0 , \EQN bc_B_ini $$
where $-\tau^{-1}_0 = B_1 / R_0 > 0$ by government
budget constraint \Ep{Gov_bc_B;a}.



\subsection{Isomorphic equilibria with fiat money 
            or unbacked debt}\label{sec:iso_mb}%

With a simple exchange of variables we can transform the equilibria with
valued fiat money in section \use{sec:Samuelson101} into equilibria
with valued unbacked government debt and vice versa. To go from a
monetary equilibrium to an equilibrium with unbacked debt, set
$R_t=p_t/p_{t+1}$ and $B_{t+1}/R_t=M/p_t$ for $t\geq 0$. To go 
in the opposite direction from an equilibrium with unbacked debt
to a monetary equilibrium, pick some arbitrary nominal quantity of 
money $M$ and solve for $p_0$ from $M/p_0=B_1/R_0$ ($= -\tau^{-1}_0$). 
All other prices, $\{p_t\}_{t\geq 1}$, are then recursively set
according to $p_{t-1}/p_t=R_{t-1}$ for $t\geq 1$.

To establish that each mapping generates an equilibrium, we need 
to show that the resulting rates of return and amounts of assets 
are the same as in the stipulated original equilibrium, as well 
as that government budget constraints are satisfied. Note that 
both mappings explicitly impose same rates of return, i.e., we
set $R_t$ equal to $p_t/p_{t+1}$ or vice versa. Furthermore, 
when going from a fiat-money equilibrium 
to an unbacked-debt equilibrium, we also set the sequences
of assets to be the same, $B_{t+1}/R_t=M/p_t$ for $t\geq 0$. So
it only remains to verify that those levels of government debt
are consistent with government budget constraints \Ep{Gov_bc_B}. 
This is trivially true for constraint \Ep{Gov_bc_B;a} in 
period $0$, given that budget constraint \Ep{Gov_bc_M} in the 
stipulated original fiat-money equilibrium holds. Regarding the 
remaining constraints \Ep{Gov_bc_B;b}, we substitute our 
mapping into the left side and the right side of each one of 
those constraints:
$${B_{t+1} \over R_t} = B_t \;\; 
\left(= R_{t-1} {B_t \over R_{t-1}} \right)
\quad \Longrightarrow \quad     
{M \over p_t} = {p_{t-1} \over p_t} {M \over p_{t-1}} ,
%%\;\; \left( = {M \over p_t}\right), 
\quad \forall t \geq 1;
$$
which confirms that the left side equals the right side for
each one of constraints \Ep{Gov_bc_B;b} and, hence, all 
government budget constraints are satisfied in the unbacked-debt
equilibrium. Thus, we conclude that an equilibrium with
fiat money can be mapped into an equilibrium with unbacked 
government debt. 

Concerning the mapping from an unbacked-debt 
equilibrium to a fiat-money equilibrium, besides having
explicitly set rates of return equal, our described 
setting of $p_0$ ensures that the government's only
budget constraint \Ep{Gov_bc_M} is satisfied;
$M/p_0=B_1/R_0 = -\tau^{-1}_0$. So it only
remains to verify that all future assets in the
monetary economy are equal to those of the
unbacked-debt equilibrium; that is, to verify that
real money supplies, $\{M/p_t\}_{t\geq 1}$, are equal 
to government indebtedness in the unbacked-debt 
equilibrium, as given by expression \Ep{Gov_debt}. 
We confirm that to be true 
by substituting our mapping into the right side of 
expression \Ep{Gov_debt} for $T \geq 1$:
$$ {B_{T+1} \over R_{T} } = -\tau^{-1}_0 \, \prod_{j=0}^{T-1} R_j 
= {M \over p_0} \, \prod_{j=0}^{T-1} {p_j \over p_{j+1} }
= {M \over p_T }.
$$
Thus, we conclude that an equilibrium with unbacked 
government debt can be mapped into an equilibrium
with fiat money. 

The isomorphic equilibria can be illustrated with our 
example of logarithmic preferences. As noted earlier,
the multiple monetary equilibria implied by
equation \Ep{samsol1} can be indexed by
$\kappa \in [0,\infty)$. The
stationary monetary equilibrium, $\kappa = 0$, has
a constant price level, i.e., a gross real rate of
return on fiat money equal to one. When mapped into
an economy with unbacked government debt, the same
stationary equilibrium allocation can be attained by 
the government at time $0$ issuing a particular 
amount of real bonds, let say, $B_1=B_H$,
as determined by the young generation's demand for
savings at a gross real interest rate equal to 
$R_0=1$. At that interest rate, the debt can be 
rolled over in perpetuity while remaining constant. As we
know from before, this high-interest-rate stationary 
equilibrium is not stable, which is also self-evident here.
If the debt issued at time $0$ is chosen to be lower 
than $B_H$, the net real interest rate will 
be negative, $R_t<1$, and, hence, subsequent debt 
levels in expression \Ep{Gov_debt} would be monotonically 
declining and converging to zero; precisely as in a monetary 
equilibrium indexed by $\kappa > 0$, where inflation 
would cause real money supplies $M/p_t$ to decline over 
time and converge to zero.

While the fiat-money equilibria and unbacked-debt 
equilibria are mathematically isomorphic, there is an
economically substantive difference. What does the
government control? In an unbacked-debt equilibrium,
the government chooses at time $0$ the quantity of 
indexed bonds $B_1$ to be issued. Furthermore, the price 
$1/R_0$ at which those bonds will be sold can be deduced
from agents' savings functions, so the government has 
full control over the size 
of the lump-sum transfer to be handed over to the 
initial old agents, $-\tau^{-1}_0 = B_1 / R_0$. 
In contrast, the government in a fiat-money 
equilibrium chooses the amount of nominal money $M$ 
to be given to the initial old agents. 
But recall from our mapping of an unbacked-debt 
equilibrium into a fiat-money equilibrium, we
started by saying: `pick some arbitrary nominal 
quantity of money $M$'; that is, the quantity $M$ itself
bears no significance. Instead, what matters is 
agents' beliefs about price levels and, hence, 
what will be the real return on money holdings.
According to equation \Ep{samsol1}, the set of 
possible price sequences consistent with money market 
clearing can be indexed by $\kappa \in [0,\infty)$. 
Thus, dependent on what agents' beliefs turn out to
be, the associated price level $p_0$ pins down the 
real value of the fiat money handed over to the 
initial old agents, $-\tau^{-1}_0 = M / p_0$. 
In that sense, the government has no direct control
over the real value of that transfer.\NFootnote{A
fiscal theory of the price level (FTPL) has been
proposed for selecting which monetary 
equilibrium that will prevail, to be expounded upon
in chapters \use{optax} and \use{fiscalmonetary}.
In the current context, FTPL would point to      
government budget constraint \Ep{Gov_bc_M}
as the determinant of price level $p_0$ and, 
consequently, of the value of $\kappa$ 
in equation \Ep{samsol1}. That is, the fiscal policy
-- here in the form of a given real transfer $-\tau^{-1}_0$
-- determines the price level sequence $\{p_t\}_{t\geq 0}$.
\index{fiscal theory of price level}}

 
%And if we include the end 
%point $\kappa=\infty$, it would be an economy where 
%money is deemed worthless and the equilibrium
%allocation is autarky.


An unbacked-debt equilibrium with declining 
indebtedness is reminiscent of a situation where a
government pays down (amortizes) its debt 
over time. Here, the size of that `amortization' 
between period $t$ and $t+1$ is
$$ {B_t \over R_{t-1}} - {B_{t+1} \over R_t}
= {B_t \over R_{t-1}} - B_t
= (1-R_{t-1})\,{B_t \over R_{t-1}} \geq 0,  \EQN amort
$$
where we have invoked equation \Ep{Gov_bc_B;b}.
Next, we will examine how the `amortization' 
is slowed down when the government makes unfunded 
purchases of goods $\{g_t\}_{t\geq 0}$, 
which amounts to deficit finance with unbacked debt. 
The analogue to such deficit finance in an isomorphic 
fiat-money equilibrium is that money is printed to 
pay for the purchases. This source of
government funding is called seigniorage and it is 
effectively a taxation of agents' money holdings,
as we will also explore below.



\section{Deficit finance with unbacked debt}

Besides the transfer to the initial old agents $-\tau^{-1}_0\geq 0$,
we now assume that the government purchases goods $g_t\geq 0$ in
period $t\geq 0$. The government budget constraints are 
$$\EQNalign{ {B_1 \over R_0}&= g_0 - \tau^{-1}_0,    \EQN Gov_bc_Bdef;a \cr
        {B_{t+1} \over R_t} &= g_t + B_t,  \quad \forall t \geq 1. 
                                                  \EQN Gov_bc_Bdef;b \cr} $$
Similar to computation \Ep{Gov_bc_iter}, we solve equations \Ep{Gov_bc_Bdef} 
forward to some date $T\geq 1$:
$$ g_0-\tau^{-1}_0 = {B_1 \over R_0}
%%={1 \over R_0}\, \left(-g_1 + {B_2 \over R_1}\right) = \ldots 
= - \sum_{i=1}^T g_i {1 \over \prod_{j=0}^{i-1} R_j} 
+  {1 \over \prod_{j=0}^{T-1} R_j}\, {B_{T+1} \over R_{T} } \EQN Gov_bc_iterdef
$$
and, hence, government indebtedness in period $T\geq 1$ equals
$$ {B_{T+1} \over R_{T} } = g_T + \sum_{i=1}^{T-1} g_i \prod_{j=i}^{T-1} R_j 
+ \left(g_0-\tau^{-1}_0\right) \, \prod_{j=0}^{T-1} R_j , \EQN Gov_debtdef
$$
which arises from the initial borrowing at time 0, $B_1/R_0=g_0-\tau^{-1}_0$,
subsequent government expenditures $\{g_t\}_{t\geq 1}^{T}$, and compounded 
interest.

The addition of government purchases does not change the 
formulation of agents' optimization problems. Since government 
purchases do not give utility, agents' objective functions are
unchanged. Budget sets are also unaffected: 
the budget constraint of an initial old is given by 
\Ep{bc_B_ini} and, at each date $t\geq 0$, a young agent faces 
budget constraints \Ep{bc_B}. Thus, an initial old agent
continues to choose $c^{-1}_0=y^{-1}_0 - \tau^{-1}_0$ and
optimal decisions by an agent born at time $t\geq 0$ are summarized
with the savings function in \Ep{savfn}:
$$ y_t^t - c_t^t = s(\alpha_t; y_t^t, y_{t+1}^t),$$
where $\alpha_t=1/R_t$. Against that backdrop, a succinct 
definition of an equilibrium is as follows.

\medskip
\noindent{\sc Definition:}  Given a sequence 
$\{g_t\}_{t=0}^\infty$ and a transfer 
$-\tau^{-1}_0\geq 0$, an equilibrium with bond-financed 
government deficits is a sequence 
$\{B_{t+1}, R_t\}_{t=0}^\infty$ such that 
(i) given $R_t$,
$$ B_{t+1} = \argmax_{b} [ u(y_t^t - b/R_t)
    + u(y^t_{t+1} + b)],  \quad \forall t \geq 0; \EQN DEq1$$
and (ii) the government budget constraints \Ep{Gov_bc_Bdef}
are satisfied.

\medskip

\noindent
Consumption of agents born in period $t\geq 0$ can be
inferred from the sequence $\{B_{t+1}, R_t\}_{t=0}^\infty$ 
and consumption of the initial old is augmented by the
transfer $-\tau^{-1}_0$. The variable $B_{t+1}$ in the definition
denotes both the solution to a representative agent's savings 
problem \Ep{DEq1} and the quantity of bonds in government 
budget constraints \Ep{Gov_bc_Bdef}, which implies that the 
bond market clears. Then, by Walras Law, the goods market also 
clears:
$$c_t^t + c_t^{t-1} + g_t = y_t^t + y_t^{t-1}, 
\quad \forall t \geq 0.        \EQN DEq3
$$


To elucidate the equilibrium dynamics of the sequence
$\{B_{t+1}, R_t\}_{t=0}^\infty$, we assume stationary
endowments and government expenditures: $y^t_t = w_1$,
$y^{t-1}_t=w_2<w_1$, and $g_t=g$ for $t\geq 0$. 
For notational simplicity we define a savings function
$f(R_t)= s(1/R_t; w_1, w_2)$. In an equilibrium, the 
demand for savings $f(R_t)$ 
equals the value of one-period government bonds issued
$B_{t+1}/R_t$. Substituting that equilibrium condition
$B_{t+1}=R_t f(R_t)$ into government budget constraints
\Ep{Gov_bc_Bdef}, we obtain the following equations that
fully characterize an equilibrium with bond-financed 
government deficits:
$$\EQNalign{f(R_0) &= g - \tau^{-1}_0,   \EQN DEq2g;a  \cr
            f(R_t) &= g + R_{t-1} f(R_{t-1}).  \quad \forall t \geq 1. 
                                          \EQN DEq2g;b \cr} $$
Next, in the spirit of our
earlier study of time 0 trading equilibria in 
section \use{sec:time0trading}, we first analyze 
stationary equilibria and then use the offer-curve approach 
of David Gale (1973) to compute nonstationary equilibria.



\subsection{Stationary equilibria and the Laffer curve}

Suppose that the economy has converged to a steady state. 
If so, equation \Ep{DEq2g;b} would have a stationary solution,
an outcome rendered reasonable by the facts that government 
purchases are constant and that $f(R_t)$ is time
invariant (because endowments are time-invariant).     
Under the supposition that $R_t=R_{t-1}=R$,
equation \Ep{DEq2g;b} becomes
$$ (1-R)\, f(R) = g.        \EQN DEq2gSS $$
This equation lends credence to earlier simile that 
declining indebtedness in an unbacked-debt equilibrium 
could be likened to a government `amortizing' its
debt by the amount $(1-R_{t-1})B_t/R_{t-1}$ in 
expression \Ep{amort}. We hinted that the government 
could instead use these `amortizations' to make 
purchases of goods. In our present stationary 
equilibrium characterized by equation \Ep{DEq2gSS},
indebtedness stays constant and, hence, the earlier
so-called `amortization' on the left side is now
fully used to purchase goods $g$ on the right side.
For later use, equation \Ep{DEq2gSS} yields a constant 
indebtedness equal to 
$$ {B \over R} = f(R) = {g \over 1-R}.    \EQN DEq2gSSd $$

For example, suppose that $u(c) = \ln(c)$.  Then
$ f(R) = {w_1 \over 2} - {w_2 \over 2 R}$.  
We have graphed $(1-R) f(R)$ in Figure \Fg{graph52f}. 
Notice that if there is one solution of equation \Ep{DEq2gSS},  
then there are at least two (except at the very peak of
the curve $(1-R) f(R)$).
When setting $g=0$ in equation \Ep{DEq2gSS}, we recover  
the pair of stationary equilibria that we have studied before. 
Specifically, in order for $(1-R) f(R)=0$, it must be true that 
either $R=1$ or $f(R)=0$. In the latter case of $f(R)=0$, we 
are evaluating the savings function at an interest $R=R_L$ at 
which agents choose to consume their own endowments, i.e., our 
earlier low-interest-rate stationary equilibrium with a relative
price of a good next period equal to $\alpha_L=1/R_L$. 
In the other case of $R=1$, it is our earlier 
high-interest-rate stationary equilibrium that we found
to be Pareto optimal and could be supported by either fiat
money or unbacked government debt. As evident from 
equation \Ep{DEq2gSS}, this socially optimal interest rate
$R=1$ can no longer be supported in an equilibrium with
bond-financed government deficits. 

\midfigure{graph52f}
\centerline{\epsfxsize=3truein\epsffile{graph52new.eps}}
\caption{The Laffer curve in revenues from the inflation tax.}
\infiglist{graph52f}
\endfigure


Given government purchases $g>0$ in 
Figure \Fg{graph52f}, the interest rates in the two possible 
stationary equilibria are $R=\underline R$ and $R=\overline R$.
While both stationary equilibria divert the same amount of 
goods $g$ from private consumption, the equilibrium with the 
higher interest rate $\overline R$ inflicts less of a distortion on
the intertemporal price. As we will further elaborate upon 
below in an isomorphic fiat-money equilibrium, $(1-R)$ can be 
interpreted as a `tax rate' and $f(R)$ as a `tax base'; thus, 
$(1-R) f(R)$ becomes a \idx{Laffer curve} showing `tax revenues'
in different stationary equilibria. Here the `tax rate' 
increases when moving from right to left on the horizontal axis 
in Figure \Fg{graph52f}, so 
$\overline R$ is on the good side of the Laffer curve and 
$\underline R$ is on the bad side.

Keep in mind that Figure \Fg{graph52f} depicts stationary 
equilibria and says nothing about how these stationary
equilibria were reached. So let's turn to the analysis
of nonstationary equilibria.



\subsection{Nonstationary equilibria}

To analyze nonstationary equilibria we adopt the offer-curve 
approach in subsection \use{sec:olg_nonstat} but instead of
a clearinghouse, we retain our current Walrasian auctioneer
who looks for an equilibrium price vector such that market 
clearing conditions \Ep{DEq3} are satisfied. Hence, in the
present context, the feasibility line 
becomes\NFootnote{Besides subtracting $g$ from the right
side of expression \Ep{Cfrontg}, we have added an extra 
initial feasibility line as compared to those in 
expression \Ep{equilibr1;b} in 
subsection \use{sec:olg_nonstat}; namely, we now require that
the consumption of the initial old and young agents in 
period $0$ satisfy feasibility with equality, i.e., the goods
market clears in period $0$.}
$$c_{t+1}^{t+1} + c_{t+1}^t 
     = w_1 + w_2 -g , \quad \forall t \geq -1,   \EQN Cfrontg
$$
as depicted in Figure \Fg{graph42f}. Notice that the endowment point 
$(y^t_t, y^t_{t+1})$ lies {\it outside} the feasibility line.
The feasibility line 
is the only object that changes as compared to earlier 
Figure \Fg{graph12f} in subsection \use{sec:olg_nonstat}. In 
particular, the offer curve remains the same for agents $t\geq 0$ 
and solves equations \Ep{offer1} for different values of
the relative price $\alpha_t = 1/ R_t$. The two intersections
of the offer curve and the feasibility line in 
Figure \Fg{graph42f} constitute the two stationary equilibria 
with interest rate $\overline R$ and $\underline R$, respectively,
as characterized above.


\midfigure{graph42f}
\centerline{\epsfxsize=3truein\epsffile{graph42.eps}}
\caption{Equilibria with debt- or money-financed government deficit finance.}
\infiglist{graph42f}
\endfigure

The steps laid out to solve for an equilibrium
allocation in subsection \use{sec:olg_nonstat} remain
the same except for the very first period $t=0$.
Instead of treating $c^0_0$ as an initial condition,
we deduce $c^0_0$ from market outcomes initiated
by the government. Specifically, as conveyed by  
equilibrium expression \Ep{DEq2g;a},  
the government borrows to finance purchases of 
goods $g>0$ and a transfer to the initial old  
$-\tau^{-1}_0\geq 0$ on the right side, which equals
the demand for savings by the young $f(R_0)$
on the left side. Thus, 
$c^0_0 = w_1 -g + \tau^{-1}_0$ ($=w_1 - B_1/R_0$)
while the consumption of the initial old
is $c^{-1}_0 = w_2 - \tau^{-1}_0$.

As an illustration, suppose that $\tau^{-1}_0=0$ and,
hence, $c^0_0 = w_1 -g$. Given that initialization 
of $c^0_0$, we can conduct a graphical analysis
in Figure \Fg{graph42f} similar to what we did
in Figure \Fg{graph22f} 
in subsection \use{sec:olg_time0log}. We will then
discover that the consumption allocation converges 
to the stationary equilibrium with interest rate 
$\underline R$. The same convergence will occur as we
successively raise the initial transfer $- \tau^{-1}_0$
until it reaches the following critical value at which 
government indebtedness in period $0$ becomes equal 
to that of the stationary equilibrium with the high 
interest rate $\overline R$:
$$ - \tau^{-1}_0 =  f(\overline R) - g 
           =  {g \over 1-\overline R} -g
           = {\overline R g \over 1- \overline R},    \EQN tau_highR
$$
where we have invoked expression \Ep{DEq2gSSd}
for the level of government indebtedness in
a stationary equilibrium. At the unique value of the
transfer in expression \Ep{tau_highR}, the economy
starts and then remains in the stationary
equilibrium with the high interest rate $\overline R$. 


In summary, there is a continuum of nonstationary solutions
of  difference equation \Ep{DEq2g}, as indexed by the
initial transfer 
$- \tau^{-1}_0 \in [0, \overline R g /(1- \overline R)]$, 
all but one of which  have
$R_t \rightarrow \underline R$ as $t \rightarrow \infty$.  
Thus, the ``bad'' Laffer-curve equilibrium is stable.

Because net interest rates are negative, government indebtedness
in any period $T\geq 1$ remains bounded in expression \Ep{Gov_debtdef}, 
where `weights' on past expenditures decline 
geometrically with the time that has passed. This also means 
that the `weight' on initial borrowing in period $0$, 
%$B_1/R_0=g_0-\tau^{-1}_0$, 
$g_0-\tau^{-1}_0$, will converge to zero. However,
as we have learned, a permanent impact of the initial 
transfer $-\tau^{-1}_0$ is if
it sets the economy on the path toward a stationary
equilibrium with a low or a high interest rate.

Solving the government budget constraint
forward in expression \Ep{Gov_bc_iterdef} yields an
unbounded present value of future expenditures 
since discounting with negative net interest rates magnify 
rather than dampen present values of future expenditures.
So what makes the present value budget constraint hold? 
It is that the bounded government indebtedness at any future
date, as just described, becomes unbounded in
its present value. Hence, the negative present value
of the stream of future government expenditures
and the positive present value of outstanding unbacked 
government debt in a distant future, both go to infinity 
in a proper way for the budget constraint to be satisfied.



\section{Deficit finance with fiat money}

Here we explore deficit finance with fiat money instead of unbacked
debt as in the preceding section. The government budget constraints 
are then
$$\EQNalign{ {M_0 \over p_0}&= g_0 - \tau^{-1}_0, \EQN Gov_bc_Mdef;a \cr
        {M_t - M_{t-1} \over p_t} &= g_t ,  \quad \forall t \geq 1. 
                                                  \EQN Gov_bc_Mdef;b \cr} 
$$
The equilibrium condition in the money market is
$f(R_t) = M_t/p_t$, which can be written
as $f(R_t) = M_{t-1} /p_t + (M_t - M_{t-1})/p_t$.  The left side
is the saving of the young.  The  term $M_{t-1} /p_t$ on  the right side is the
dissaving of the old (the real value of currency  that they
exchange for time $t$ consumption). The  term $(M_t - M_{t-1})/p_t$  on the right
is the dissaving of the government (its deficit), which is
the real value of the additional currency that the government
prints at $t$ and uses to purchase time $t$ goods from the young.

\medskip
\noindent{\sc Definition:}  Given a sequence 
$\{g_t\}_{t=0}^\infty$ and a transfer 
$-\tau^{-1}_0\geq 0$, an equilibrium with money-financed 
government deficits is a sequence 
$\{M_t, p_t\}_{t=0}^\infty$
with $0 < p_t < +\infty$ and $M_t >0$ such that 
(i) given $\{p_t, p_{t+1}\}$,
$$ M_t = \argmax_{m \geq 0} \left[ u(y_t^t - m/p_t)
       + u(y^t_{t+1} + m/p_{t+1})\right],
       \quad \forall t \geq 0; \EQN equilib1;a $$
and (ii) the government budget constraints \Ep{Gov_bc_Mdef}
are satisfied.

\medskip


\noindent
Notice that money neutrality prevails in this economy. That is, 
if there exists an equilibrium for some sequence
$\{M_t, p_t\}_{t=0}^\infty$ then there are alternative sequences 
$\{\kappa M_t, \kappa p_t\}_{t=0}^\infty$, for any $\kappa\in (0, \infty)$,
which also constitute equilibria with the same consumption allocation.
This follows from the fact that the variables for nominal money
supply and price level show up as ratios of one another in all
relevant instances.



\subsection{Isomorphic equilibria of deficit finance}

As we might anticipate from subsection \use{sec:iso_mb}, there
are isomorphic equilibria of deficit finance where the 
government uses either fiat money or unbacked debt. That is, 
the government finances deficits either by printing money or by
borrowing at a negative real net interest rate. 

To go from an equilibrium with unbacked debt to an equilibrium 
with fiat money, pick some arbitrary nominal quantity of 
money $M_0$ and solve for $p_0$ from $M_0/p_0=B_1/R_0$ ($= g_0-\tau^{-1}_0$). 
All other prices, $\{p_t\}_{t\geq 1}$, are then recursively set
according to $p_{t-1}/p_t=R_{t-1}$ for $t\geq 1$. Thereafter,
nominal money supplies are set to satisfy 
$M_t / p_t = B_{t+1} / R_t$.

While our mapping into a fiat-money equilibrium imposes the same 
rates of return and amounts of assets as in the original 
unbacked-debt equilibrium, we still have to establish that
government budget constraints \Ep{Gov_bc_Mdef} are satisfied. 
This is trivially true for constraint \Ep{Gov_bc_Mdef;a} in 
period $0$, given our above setting of $p_0$. Regarding the 
remaining constraints \Ep{Gov_bc_Mdef;b}, we substitute our 
mapping into the left side of those budget constraints
(to confirm that this generates the right side):
$${M_t - M_{t-1} \over p_t} = {M_t \over p_t} 
                            - {M_{t-1} \over p_t} {p_{t-1} \over p_{t-1}}
                            = {B_{t+1} \over R_t} 
                            - {B_t \over R_{t-1}}\,R_{t-1} 
                            = g_t 
$$
where the last equality invokes budget constraints 
\Ep{Gov_bc_Bdef;b} in the original unbacked-debt equilibrium.
Thus, we conclude that an equilibrium with bond-financed
government deficits can be mapped into an equilibrium with
fiat money replacing the unbacked debt. 

We invite the
reader to do the mapping in the opposite direction
from an equilibrium with money-financed government
deficits to one with unbacked government debt.


\subsection{Computing equilibria}

To compute an equilibrium with a constant stream of
government expenditures $g$, write equation \Ep{Gov_bc_Mdef;b} as
$$ {M_t \over  p_t} = {M_{t-1} \over p_{t-1} }  {p_{t-1} \over p_t}
     + g,  \quad \forall t \geq 1;$$
and ${M_0 / p_0}$ is given by equation \Ep{Gov_bc_Mdef;a}.
Substitute the equilibrium condition $M_t / p_t = f(R_t)$, where
$R_t=p_t/p_{t+1}$, into these equations to
get
$$\EQNalign{ f(R_t) &= f(R_{t-1})R_{t-1} + g  \quad 
                             \forall t \geq 1\EQN money1;a \cr
             f(R_0) &=  g - \tau^{-1}_0. \EQN money1;b \cr}
$$
As one should have anticipated, the equilibrium characterization
in equations \Ep{money1} is identical
to that of an equilibrium with bond-financed deficits
in equations \Ep{DEq2g}.

\medskip
\noindent
{\bf [Elaborate]}
\medskip

Next, we revisit the example of logarithmic preferences in
subsection \use{sec:Meqlog}. Recall difference equation 
\Ep{samdiff1} that we here formulate with a lag operator
$$\left( 1 - {w_2 \over w_1} L^{-1} \right) p_t = {2 \over w_1} M_t,
                                                       \EQN samdiff2
$$
where we also allow for our present case of a time-varying forcing term
$M_t$. The solution to this difference equation is
$$ p_t = {2 \over w_1} \sum_{i=0}^\infty \left({w_2 \over w_1}\right)^i M_{t+i}
      + \kappa \left({w_1\over w_2}\right)^t, \EQN samsol2 $$
If the nominal money supply stays constant, the solution simplifies
to earlier expression \Ep{samsol1}, but otherwise we need the infinite sequence
of nominal money supplies. To illustrate such an equilibrium
calculation we proceed under the assumption that the economy has converged 
to a stationary equilibrium. Since the rate of return (the
inverse of the inflation rate) stays constant in a stationary equilibrium,
$p_t/p_{t+1}=R$, it is natural to conjecture a constant growth rate of the
nominal money supply, $M_{t+1} = zM_t$ for some $z>0$. According to
the first proposition in section \use{sec:olg_optimality}, a necessary 
and sufficient condition for the existence of a monetary equilibrium
is then that $z u'(y^i_i)/u'(y^i_{i+1})<1$, which here translates
to  $z w_2/w_1<1$.
Given the sequence $\{M_{t+i}\}_{i=0}^\infty=\{z^i M_t\}_{i=0}^\infty$
where $z<w_1/w_2$,  we solve equation \Ep{samsol2} to obtain  
price sequences consistent with money market clearing,
$$ p_t = {2 \over w_1} { M_t \over 1 - {w_2 \over w_1} z }
      + \kappa \left({w_1\over w_2}\right)^t. \EQN samsol2z 
$$
\medskip
\noindent
{\bf [Elaborate on this. Then connect back to graphs \Fg{graph52f}
and \Fg{graph42f}. Finally, circle back to that the ``bad'' Laffer curve 
equilibrium is stable under the ``perfect foresight'' dynamics.]}


The undesirable stability of the bad Laffer curve equilibrium is a feature of
perfect foresight dynamics.
Bruno and Fischer (1990) and Marcet and Sargent (1989) analyze
how the system behaves under two different types of adaptive
dynamics.  They find that either under a crude form
of  adaptive expectations  or under a least-squares learning scheme,
$R_t$ converges to  $\overline R$.   This finding  is comforting because
the comparative dynamics are more plausible at $\overline  R$
(larger deficits bring higher inflation).  Furthermore,
Marimon and Sunder (1993) present experimental evidence
pointing toward the selection  made by the adaptive dynamics.
Marcet and Nicolini  (2003) build and calibrate an adaptive model of several
Latin American hyperinflations that rests on this selection. Sargent, Williams, and Zha (2009)
extend and estimate that model.
\auth{Williams, Noah}%
\auth{Zha, Tao}%
\auth{Marcet, Albert}
\auth{Nicolini, Juan Pablo}
\auth{Marimon, Ramon}
\auth{Sunder, Shyam}
\auth{Bruno, Michael}
\auth{Fischer, Stanley}
\auth{Sargent, Thomas J.}








%\section{Pay-as-you-go social security}









\section{Remain to be revised}

\index{Lucas tree!in overlapping generations model}%
\subsection{Another Gale diagram}\label{sec:olg_lucas_tree}%
To set up an environment with a ``Lucas tree,'' take the preference and endowment
structure to be the same as example 1 and modify only one feature.
Endow the initial old with a claim
to a constant stream of $d >0$ units of consumption for each
$t\geq 1$.\NFootnote{This is a version of an example of Brock (1990). The `Lucas tree' refers to a colorful interpretation
of a dividend stream as `fruit' falling from a `tree' in a pure exchange economy studied by Lucas (1978). See chapter \use{assetpricing1}.}
\auth{Brock, William A.}%
\auth{Lucas, Robert E., Jr.}%
   Thus, the budget constraint of an initial old person now becomes
$$ q_1^0 c_1^0 = d \sum_{t=1}^\infty q_t^0 + q_1^0 y_1^0. $$
The offer curve of each young agent remains as before,
but now the feasibility line is
$$ c_i^i + c^{i-1}_i =  y_o + y_y + d $$
for all $i \geq 1$.  Note that young agents are  endowed
below the feasibility line in  Figure \Fg{graph32f}, %Figure 8.3,
It seems that there
are  two candidates for stationary equilibria, one
with constant  $\alpha < 1$, another with constant $\alpha >1$.
The one with $\alpha  <1$ is associated with the
steeper budget line in Figure \Fg{graph32f}. %Figure 8.3.
 However, the candidate stationary
equilibrium with $ \alpha >1 $ cannot be an equilibrium. % for a reason
%similar to that encountered in example 3. 
 At the price
system associated with an $\alpha >1$, the wealth of the initial
old would be unbounded, which would prompt them to consume an
unbounded amount, which is not feasible.
This argument rules out not only the stationary $\alpha >1$ equilibrium
but also all nonstationary candidate equilibria that converge to
that constant $\alpha$.  Therefore, there is a unique equilibrium;
it is stationary and has $\alpha < 1$.

%%%%%%%
%\topinsert{
%$$
%\grafone{graph32.eps,height=2.8in}{{\bf Figure 8.3} Unique equilibrium with
%a fixed-dividend asset.}
%$$
%}\endinsert

\subsection{Endowment at $+\infty$}\label{sec:endowment_infty}%
To get  into the mood to think about some issues to be discussed in
 section \use{sec:olg_optimality}, take the preference and endowment structure
of our  example 1 above and modify only one feature. Change the
endowment of the initial old to be  $y^0_1 =\epsilon>0 $ {\it and}
``$\delta=1-\epsilon >0$ units of consumption at $t=+\infty$,''
by which we mean that we take
$$ \sum_t q_t^0 y_t^0 = q_1^0 \epsilon + \delta \lim_{t\rightarrow
   \infty} q_t^0.$$
It is easy to verify that the only competitive equilibrium
of the economy with this specification of endowments
has $q_t^0 = 1 \ \forall t\geq 1$, and thus $\alpha_t = 1
\ \forall t \geq 1$.   This situation reflects that
  all  of the ``low-interest-rate'' equilibria  that we computed in example 1
would assign an infinite value to the endowment
of the initial old.  Confronted with such prices, the
initial old would demand unbounded consumption, which would not be 
feasible.  Therefore, that cannot be an equilibrium.

\midfigure{graph32f}
\centerline{\epsfxsize=3truein\epsffile{graph32.eps}}
\caption{Unique equilibrium with a fixed-dividend asset.}
\infiglist{graph32f}
\endfigure

    If we interpret the gross rate of return on the
tree as $\alpha^{-1} = {p + d  \over p}$, where $p = \sum_{t=1}^\infty
q_t^0 d$, we can compute that $p = {d \over R-1}$ where
$R = \alpha^{-1}$.  Here $p$ is the price of the
Lucas tree for all $t \geq 1$.

    

\subsection{Growth}

  It is easy to extend these models to the case  in which  there is
growth in the population.  Let there be $N_t = n N_{t-1}$ identical young
people at time $t$, with $n > 0$.   For example, consider the economy
with money-financed deficits.  The total money supply
is $N_t M_t$, and the government budget constraint is
$$ N_t M_t - N_{t-1} M_{t-1} = N_t p_t g, $$
where $g$ is per-young-person government   purchases.
Dividing both sides of the budget constraint by
$ N_t$ and rearranging gives
$$ {M_t \over p_{t+1}} { p_{t+1} \over p_t} = n^{-1} {M_{t-1} \over p_t}
     + g .\EQN equilib1c  $$
This equation replaces equation \Ep{equilib1;b} in the definition of
an equilibrium with money-financed deficits.  (Note that
in a steady state, $R=n$ is the high-interest-rate equilibrium.)
Similarly, in the economy with bond-financed deficits,
the government budget constraint would become
$$ {B_{t+1} \over R_t} = n^{-1} B_t + g_t .$$

  It is also easy to modify things to permit the government to
tax young and old people at $t$.  In that case, with government
bond finance the government budget constraint
becomes
$$ {B_{t+1} \over R_t} = n^{-1} B_t + g_t - \tau^t_t - n^{-1} \tau^{t-1}_t ,$$
where $\tau^s_t$ is the time $t$ tax on a person born in period $s$.

\auth{Wallace, Neil}%

\section{Optimality and  existence of monetary equilibria}\label{sec:olg_optimality}%
Wallace (1980)
discusses a connection between nonoptimality of
the equilibrium without valued money and the existence of monetary
equilibria. For an economy not having the  storage technology available in  Wallace's economy,
we study how  arguments similar to his apply. The
environment is as follows. At any
date $t$, the population consists of $N_t$ young agents and
$N_{t-1}$ old agents where $N_t=n N_{t-1}$ with $n > 0$. Each young
person is endowed with $y_1> 0$ goods, and an old person receives the
endowment $y_2> 0$. Preferences of a young agent at time $t$ are given
by the utility function $u(c^t_t,c^t_{t+1})$, which is twice
differentiable with indifference curves that are convex to the
origin. The two goods in the utility function are
normal goods, and $$\theta(c_1,c_2) \equiv u_1(c_1,c_2)/u_2(c_1,c_2),$$
the marginal rate of substitution function, approaches infinity as
$c_2/c_1$ approaches infinity and approaches zero as $c_2/c_1$ approaches
zero. The welfare of the initial old agents at time $1$ is
strictly increasing in $c^0_1$, and each one of them is endowed with
$y_2$ goods and $m^0_0> 0$ units of fiat money. Thus, the
beginning-of-period aggregate nominal money balances in the initial
period $1$ are $M_0=N_0 m^0_0$.

For all $t \geq 1$, $M_t$, the post-transfer time $t$ stock of fiat
money obeys $M_t = z M_{t-1}$ with $z > 0$. The time $t$ transfer
(or tax), $(z-1) M_{t-1}$, is divided equally at time $t$ among the
$N_{t-1}$ members of the current old generation. The transfers
(or taxes) are fully anticipated and are viewed as lump-sum: they do
not depend on consumption and saving behavior. The budget constraints
of a young agent born in period $t$ are
$$ \EQNalign{ c_t^t + {m_t^t \over p_t} & \leq y_1, \EQN wall1 \cr
              c_{t+1}^t & \leq y_2 + {m_t^t \over p_{t+1} }
               + {(z-1) \over N_t} {M_t \over p_{t+1} }, \EQN wall2 \cr
              m_t^t & \geq 0,                            \EQN wall3 \cr} $$
where $p_t> 0$ is the time $t$ price level. In a nonmonetary equilibrium,
the price level is infinite, so the real values of both money holdings and
transfers are zero.
Since all members in a generation are identical, the nonmonetary
equilibrium is autarky with a marginal rate of substitution equal to
$$
\theta_{\rm aut} \equiv { u_1(y_1,y_2) \over u_2(y_1,y_2)}.
$$
We ask two questions about this economy. Under what
circumstances does a monetary equilibrium exist? And, when it exists,
under what circumstances does it improve matters?

Let $\hat m_t$ denote the equilibrium real money
balances of a young agent at time $t$,
$\hat m_t \equiv M_t / (N_t p_t)$. Substitution of equilibrium money
holdings into budget constraints \Ep{wall1} and \Ep{wall2} at equality
yield $c^t_t= y_1 - \hat m_t$ and $c^t_{t+1}=y_2 + n \hat m_{t+1}$.
In a monetary equilibrium, $\hat m_t > 0$ for all $t$ and the
marginal rate of substitution $\theta(c^t_t,c^t_{t+1})$ satisfies
$$
\theta(y_1- \hat m_t, \ y_2 + n \hat m_{t+1}) = { p_t \over p_{t+1} }
> \theta_{\rm aut},   \quad  \forall t \geq 1 .            \EQN wall5
$$
The equality part of \Ep{wall5} is the first-order condition for money
holdings of an agent born in period $t$ evaluated at the equilibrium
allocation. Since $c^t_t <y_1$ and $c^t_{t+1} > y_2$ in a monetary
equilibrium, the inequality in
\Ep{wall5} follows from the assumption that the two goods in the
utility function are normal goods.

Another useful characterization of the equilibrium rate of return on
money, $p_t/p_{t+1}$, can be obtained as follows.
By the rule generating $M_t$ and the equilibrium condition
$M_t/p_t=N_t \hat m_t$, we have for all $t$,
$$
{p_t \over p_{t+1}} = {M_{t+1} \over z M_t} {p_t \over p_{t+1}}
= {N_{t+1} \hat m_{t+1} \over z N_t \hat m_t}
= {n \over z} {\hat m_{t+1} \over \hat m_t}.                 \EQN wall6
$$
We are now ready to address our first question, under what circumstances
does a monetary equilibrium exist?

\medskip
\noindent{\sc Proposition:}  $\theta_{\rm aut} z < n$ is necessary
and sufficient for the existence of at least one monetary equilibrium.
\medskip

\noindent{\sc Proof:}
We first establish necessity. Suppose to the contrary that there is
a monetary equilibrium and $\theta_{\rm aut} z / n \geq 1$. Then,
by the inequality part of \Ep{wall5} and expression \Ep{wall6},
we have for all $t$,
$$
{\hat m_{t+1} \over \hat m_t}  > { z \theta_{\rm aut} \over n} \geq 1. \EQN wallmt
$$
If $z \theta_{\rm aut} / n > 1$, one plus the net growth rate of $\hat m_t$ is
bounded uniformly above one and, hence,
the sequence $\{\hat m_t\}$ is unbounded, which is
inconsistent with an equilibrium because real money balances per capita
cannot exceed the endowment $y_1$ of a young agent.
If $z \theta_{\rm aut} / n = 1$, the strictly increasing sequence
$\{\hat m_t\}$ in \Ep{wallmt} might not be unbounded but converge
to some constant $\hat m_\infty$.
According to \Ep{wall5} and \Ep{wall6}, the marginal rate
of substitution will then converge to $n/z$, which by assumption is now
equal to $\theta_{\rm aut}$, the marginal rate of substitution in
autarky. Thus, real balances must be zero in the limit, which contradicts
the existence of a strictly increasing sequence of positive real balances
in \Ep{wallmt}.

To show sufficiency, we prove the existence of a unique
equilibrium with constant per capita real money balances when $\theta_{\rm aut} z < n$.
Substitute our candidate equilibrium,
$\hat m_t = \hat m_{t+1} \equiv \hat m$,  into \Ep{wall5} and \Ep{wall6},
which yields two equilibrium conditions,
$$
\theta(y_1- \hat m, \ y_2 + n \hat m) = { n \over z } > \theta_{\rm aut}.
$$
The inequality part is satisfied under the parameter restriction of
the proposition, and we only have to show the existence of
$\hat m \in [0, y_1]$ that satisfies the equality part.
But the existence (and uniqueness) of such a $\hat m$ is trivial.
Note that the marginal rate of substitution on the left side of
the equality is equal to $\theta_{\rm aut}$ when $\hat m=0$.
Next, our assumptions on preferences imply that the marginal rate
of substitution is strictly increasing in $\hat m$, and approaches
infinity when $\hat m$ approaches $y_1$. \qed
\medskip

The stationary monetary equilibrium in the proof will be referred to
as the $\hat m$ equilibrium. In general, there are other nonstationary
monetary equilibria when the parameter condition of the proposition
is satisfied. For example, in the case of logarithmic preferences and
a constant population, recall the continuum of equilibria
indexed by the scalar $\kappa > 0$ in expression \Ep{samsol1}. But here we
choose to focus solely on the stationary $\hat m$ equilibrium and
its welfare implications. The $\hat m$ equilibrium will be compared
to other feasible
allocations using the Pareto criterion. Evidently, an allocation
$C = \{c^0_1; (c^t_t,c^t_{t+1}), t\geq 1\}$ is feasible if
$$ \EQNalign{
N_t c_t^t + N_{t-1} c_t^{t-1} &\leq N_t y_1 + N_{t-1} y_2,
                                       \quad  \forall t \geq 1, \cr
\noalign{\hbox{\rm or, equivalently,}}
n c_t^t + c_t^{t-1} &\leq n y_1 + y_2, \quad  \forall t \geq 1.
                                                        \EQN wall7 \cr}
$$
The definition of Pareto optimality is:

\medskip
\noindent{\sc Definition:} A feasible allocation $C$ is Pareto optimal
if there is no other feasible allocation $\tilde C$ such that
$$\EQNalign{
&\tilde c^0_1 \geq c^0_1, \cr
&u(\tilde c^t_t, \tilde c^t_{t+1}) \geq u(c^t_t,c^t_{t+1}),
                \quad  \forall t \geq 1, \cr}
$$
and at least one of these weak inequalities holds with strict inequality.
\medskip

We first examine under what circumstances the nonmonetary
equilibrium (autarky) is Pareto optimal.

\medskip
\noindent{\sc Proposition:}  $\theta_{\rm aut} \geq n$ is necessary
and sufficient for the optimality of the nonmonetary equilibrium
(autarky).
\medskip

\noindent{\sc Proof:}
To establish sufficiency, suppose to the contrary that there exists
another feasible allocation $\tilde C$ that is Pareto superior to
autarky and $\theta_{\rm aut} \geq n$. Without loss of generality,
assume that the allocation $\tilde C$ satisfies \Ep{wall7} with
equality. (Given an allocation that is Pareto superior to autarky but
that does not satisfy \Ep{wall7}, one can easily construct another
allocation that is Pareto superior to the given allocation,
and hence to autarky.) Let period $t$ be the first
period when this alternative allocation $\tilde C$ differs from the
autarkic allocation. The requirement that the old generation in this
period is not made worse off, $\tilde c^{t-1}_t \geq y_2$, implies
that the first perturbation from the autarkic allocation must be
$\tilde c^t_t < y_1$, with the subsequent implication that
$\tilde c^t_{t+1} > y_2$. It follows that the consumption of
young agents at time $t+1$ must also fall below $y_1$, and we define
$$
\epsilon_{t+1} \equiv y_1 - \tilde c^{t+1}_{t+1} > 0.      \EQN wall8
$$
Now, given $\tilde c^{t+1}_{t+1}$, we compute the smallest number
$c^{t+1}_{t+2}$ that satisfies
$$
u(\tilde c^{t+1}_{t+1},c^{t+1}_{t+2}) \geq u(y_1,y_2).
$$
Let $\overline c^{t+1}_{t+2}$ be the solution to this problem.
Since the allocation $\tilde C$ is Pareto superior to autarky,
we have $\tilde c^{t+1}_{t+2} \geq \overline c^{t+1}_{t+2}$. Before using
this inequality, though, we want to derive a convenient expression
for $\overline c^{t+1}_{t+2}$.

Consider the indifference curve of $u(c_1,c_2)$ that yields a fixed
utility equal to $u(y_1,y_2)$.
In general, along an indifference curve, $c_2=h(c_1)$,
where $h'=-u_1/u_2=-\theta$ and $h'' > 0$. Therefore, applying the
intermediate value theorem to $h$, we have
$$
h(c_1) = h(y_1) + (y_1-c_1) [-h'(y_1) + f(y_1-c_1)],       \EQN wall9
$$
where the function $f$ is strictly increasing and $f(0)=0$.

Now, since $(\tilde c^{t+1}_{t+1},\overline c^{t+1}_{t+2})$ and
$(y_1,y_2)$ are on the same indifference curve, we can use
\Ep{wall8} and \Ep{wall9} to write
$$
\overline c^{t+1}_{t+2} = y_2 + \epsilon_{t+1} [\theta_{\rm aut}
                                 + f(\epsilon_{t+1})],
$$
and after invoking $\tilde c^{t+1}_{t+2} \geq \overline c^{t+1}_{t+2}$,
we have
$$
\tilde c^{t+1}_{t+2} - y_2 \geq \epsilon_{t+1} [\theta_{\rm aut}
                                 + f(\epsilon_{t+1})].       \EQN wall10
$$
Since $\tilde C$ satisfies \Ep{wall7} at equality, we also have
$$
\epsilon_{t+2} \equiv y_1 - \tilde c^{t+2}_{t+2}
= {\tilde c^{t+1}_{t+2} - y_2 \over n}.                     \EQN wall11
$$
Substitution of \Ep{wall10} into \Ep{wall11} yields
{\ninepoint
$$\eqalign{ \epsilon_{t+2} & \geq  \epsilon_{t+1}
                 {\theta_{\rm aut} + f(\epsilon_{t+1}) \over n} \cr
                & > \epsilon_{t+1},\cr}\EQN wall12
$$
}% endninepoint
where the strict inequality follows from
$\theta_{\rm aut} \geq n$ and $f(\epsilon_{t+1})>0$ (implied by
$\epsilon_{t+1}>0$).
Continuing these computations of successive values of $\epsilon_{t+k}$ yields
{\ninepoint
$$
\epsilon_{t+k} \geq \epsilon_{t+1}
         \prod_{j=1}^{k-1}
{\theta_{\rm aut} + f(\epsilon_{t+j}) \over n}
 > \epsilon_{t+1} \left[{\theta_{\rm aut} +
 f(\epsilon_{t+1}) \over n} \right]^{k-1},
\ \hbox{\rm for}\;k>2,
$$
}%endninept
where the strict inequality follows from the fact that $\{\epsilon_{t+j}\}$
is a strictly increasing sequence.
Thus, the $\epsilon$ sequence is bounded below by a strictly increasing
exponential and hence is unbounded. But such an unbounded sequence violates
feasibility because $\epsilon$ cannot exceed the endowment $y_1$ of a
young agent. It follows that we can rule out the existence of a
Pareto superior allocation $\tilde C$, and conclude that
$\theta_{\rm aut} \geq n$ is
a sufficient condition for the optimality of autarky.

To establish necessity, we prove the existence of an alternative
feasible allocation $\hat C$ that is Pareto superior to autarky when
$\theta_{\rm aut} < n$. First, pick an $\epsilon > 0$ sufficiently
small so that
$$
\theta_{\rm aut} + f(\epsilon) \leq n,                      \EQN wall13
$$
where $f$ is defined implicitly by equation \Ep{wall9}.
Second, set $\hat c^t_t= y_1 - \epsilon \equiv \hat c_1$, and
$$
\hat c^t_{t+1} = y_2 + \epsilon [\theta_{\rm aut} + f(\epsilon)]
               \equiv \hat c_2,    \quad  \forall t \geq 1. \EQN wall14
$$
That is, we have constructed a consumption bundle
$(\hat c_1, \hat c_2)$ that lies on the same indifference curve as
$(y_1, y_2)$, and from \Ep{wall13} and \Ep{wall14}, we have
$$
\hat c_2 \leq y_2 + n \epsilon,
$$
which ensures that the condition for feasibility \Ep{wall7} is satisfied
for $t\geq 2$. By setting $\hat c^0_1 = y_2 + n \epsilon$, feasibility
is also satisfied in period $1$ and the initial old generation is
then strictly better off under the alternative allocation $\hat C$.
\qed

\medskip
With a constant nominal money supply, $z=1$, the two
propositions show that a monetary equilibrium exists if and only
if the nonmonetary equilibrium is suboptimal. In that case, the
following proposition establishes that the stationary $\hat m$
equilibrium is optimal.

\medskip
\noindent{\sc Proposition:}  Given $\theta_{\rm aut} z < n$, then
$z \leq 1$ is necessary and sufficient for the optimality of the
stationary monetary equilibrium $\hat m$.
\medskip

\noindent{\sc Proof:}
The class of feasible stationary allocations with
$(c^t_t,c^t_{t+1}) = (c_1,c_2)$ for all $t\geq 1$, is given by
$$
c_1 + {c_2 \over n} \leq y_1 + {y_2 \over n},        \EQN wall15
$$
i.e., the condition for feasibility in \Ep{wall7}. It follows that
the $\hat m$ equilibrium satisfies \Ep{wall15} at equality, and
we denote the associated consumption allocation of an agent born
at time $t\geq 1$ by $(\hat c_1, \hat c_2)$. It is also the case
that $(\hat c_1, \hat c_2)$ maximizes an agent's utility subject
to budget constraints \Ep{wall1} and \Ep{wall2}.
The consolidation of these two constraints yields
$$
c_1 + {z \over n} c_2 \leq y_1 + {z \over n} y_2
    +  {z \over n} {(z-1) \over N_t} {M_t \over p_{t+1} },
                                                     \EQN wall16
$$
where we have used the stationary rate or return in \Ep{wall6},
$p_t/p_{t+1}=n/z$. After also invoking $z M_t=M_{t+1}$,
$n=N_{t+1}/N_t$, and the equilibrium condition
$M_{t+1}/(p_{t+1} N_{t+1})= \hat m$, expression \Ep{wall16}
simplifies to
$$
c_1 + {z \over n} c_2 \leq y_1 + {z \over n} y_2 + (z-1) \hat m.
                                                      \EQN wall17
$$

To prove the statement about necessity, Figure \Fg{graph16f} %Figure 8.6
depicts the two curves
\Ep{wall15} and \Ep{wall17} when condition $z \leq 1$ fails to hold,
i.e., we assume that $z>1$. The point that maximizes utility subject
to \Ep{wall15} is denoted $(\overline c_1, \overline c_2)$.
Transitivity of preferences  and the fact that the
slope of budget line \Ep{wall17} is flatter than that of \Ep{wall15}
imply that $(\hat c_1, \hat c_2)$ lies southeast of
$(\overline c_1, \overline c_2)$. By revealed preference, then,
$(\overline c_1, \overline c_2)$ is preferred to
$(\hat c_1, \hat c_2)$ and all generations born in period $t \geq 1$
are better off under the allocation $\overline C$.
The initial old generation can also be made better off under this
alternative allocation since it is feasible to strictly
increase their consumption,
$$
\overline c^0_1 = y_2 + n(y_1 - \overline c^1_1)
                > y_2 + n(y_1 - \hat c^1_1) = \hat c^0_1.
$$
Thus, we have established that $z \leq 1$ is necessary for the
optimality of the stationary monetary equilibrium $\hat m$.

To prove sufficiency, note that \Ep{wall5}, \Ep{wall6} and $z\leq 1$
imply that
$$
\theta(\hat c_1, \hat c_2) = { n \over z} \geq n.
$$
We can then construct an argument that is analogous to the sufficiency
part of the proof to the preceding proposition. \qed

%%%%%%%%%%
%$$
%\grafone{graph16.eps,height=2.5in}{{\bf Figure 8.6} The feasibility
%line \Ep{wall15} and the budget line \Ep{wall17} when $z>1$.
%The consumption allocation in
%the monetary equilibrium is $(\hat c_1, \hat c_2)$, and the point
%that maximizes utility subject to the feasibility line
%is denoted $(\overline c_1, \overline c_2)$.}
%$
%%%%%%%%

\midfigure{graph16f}
\centerline{\epsfxsize=3truein\epsffile{graph16.eps}}
\caption{The feasibility line \Ep{wall15} and the budget line \Ep{wall17}
when $z>1$.  The consumption allocation in the monetary equilibrium is
$(\hat c_1, \hat c_2)$, and the point that maximizes utility subject to
the feasibility line is denoted $(\overline c_1, \overline c_2)$.}
\infiglist{graph16f}
\endfigure

As pointed out by Wallace (1980), the proposition implies no connection
between the path of the price level in an $\hat m$ equilibrium and the
optimality of that equilibrium.  Thus, there may be an optimal monetary
equilibrium with positive inflation, for example, if
$\theta_{\rm aut} < n < z \leq 1$; and there may be a nonoptimal
monetary equilibrium with a constant price level,  for example,
if $z = n > 1 > \theta_{\rm aut}$.  What counts is the nominal
quantity of fiat money. The proposition suggests that the quantity
of money should not be increased. In particular, if $z \leq 1$,
then an optimal $\hat m$ equilibrium exists whenever the nonmonetary
equilibrium is nonoptimal.

\auth{Balasko, Yves} \auth{Shell, Karl}\auth{Cass, David}
\subsection{Balasko-Shell criterion for optimality}
For the case of constant population, Balasko and Shell (1980) have
established a convenient general criterion for testing whether
allocations are optimal.\NFootnote{Balasko and Shell credit David
Cass (1971) with having authored a version of their criterion.}
Balasko and Shell permit diversity among agents in terms of
endowments $[w^{th}_t, w^{th}_{t+1}]$ and utility functions
$u^{th}(c^{th}_t, c^{th}_{t+1})$, where $w^{th}_s$ is the time $s$
endowment of an agent named $h$ who is born at $t$ and $c^{th}_s$
is the time $s$ consumption of agent named $h$ born at $t$.
Balasko and Shell
 assume  fixed populations of types $h$ over time. They
impose several kinds of technical
conditions that serve to rule out possible pathologies.  The two main
ones are these.
First, they assume that indifference curves have neither flat parts nor kinks,
and they also rule out indifference curves with flat parts or kinks as limits
of sequences of indifference curves for given $h$ as $t\to\infty$.  Second,
they assume that the aggregate endowments $\sum_h (w^{th}_t + w^{t-1,h}_t)$
are uniformly bounded from above and that there exists an $\epsilon>0$
 such that
$w^{sh}_t>\epsilon$ for all $s,h$, and for $t\in\{s,s+1\}$.  They consider
consumption allocations
uniformly bounded away from the axes.  With these conditions, Balasko and Shell
consider the class of allocations in which all young agents at $t$ share a
common marginal rate of substitution $1+r_t
\equiv u^{th}_1(c^{th}_t,c^{th}_{t+1})
/u^{th}_2(c^{th}_t,c^{th}_{t+1})$ and in which all of the endowments
are consumed.  Then Balasko and Shell show that an allocation is Pareto
optimal if and only if
$$\sum_{t=1}^\infty\ \prod_{s=1}^t  [1+r_s]=+\infty,
\EQN balshell$$
that is, if and only if the infinite sum of $t$-period gross interest rates,
$\prod_{s=1}^t [1+r_s]$, diverges.

The Balasko-Shell criterion for optimality succinctly summarizes the sense in
which low-interest-rate economies are not optimal.  We have already encountered
repeated examples of the situation that, before an equilibrium with valued
currency can exist, the equilibrium without valued currency must be a
low-interest-rate economy in just the sense identified by Balasko and Shell's
criterion, \Ep{balshell}.  Furthermore, by applying the
Balasko-Shell criterion,
\Ep{balshell},
 or  generalizations of it that allow for a positive net growth
rate of population $n$, it can be shown that, among equilibria with valued
currency, only equilibria with high rates of return on currency are optimal.

\auth{Smith, Bruce D.}  \auth{Wallace, Neil}
\section{Within-generation heterogeneity}\label{sec:realbills}%
This section describes an overlapping generations
model having within-generation heterogeneity of
endowments.  We shall follow Sargent and Wallace  (1982)
and Smith (1988)
and use this model as a vehicle for talking about some
issues  in monetary theory that  require  a setting
in which government-issued  currency coexists with and is a more-or-less
 good substitute for
private IOUs.

We now assume  that
within each generation born at $t  \geq 1$,
there are $J$ groups of agents. There
is a constant number $N_j$ of group $j$ agents.  Agents of group
$j$ are endowed with $w_1(j)$ when young and
$w_2(j)$ when old.  The saving function
of a consumer of group $j$ born at time $t$
solves the time $t$ version of  problem
\Ep{savingsprob}.    We denote this savings
function $f(R_t,j)$. If we  assume that all
consumers of generation $t$  have
preferences  $U^t(c^t) = \ln c_t^t + \ln c^t_{t+1}$,
the saving function is
$$ f(R_t,j) = .5\left( w_1(j) - {w_2(j)  \over R_t} \right).$$
At $t=1$, there are old people who are endowed in
the aggregate with $H = H(0)$ units of an inconvertible currency.

For example, assume that $J=2$, that $(w_1(1), w_2(1))  =
 (\alpha, 0)$ and that $
(w_1(2), w_2(2))   = (0,  \beta)$, where $\alpha >0, \beta >0$.
The type 1 people are lenders, while the type 2 are borrowers.
For the case of log preference we have the savings
functions
$f(R,1) = \alpha /2, f(R,2) = - \beta/(2R)$.

\subsection{Nonmonetary equilibrium}

A nonmonetary equilibrium consists of sequences $(R, s_j)$ of rates of return $R$
and savings rates for  $j= 1, \ldots , J$  and $t \geq 1$ that satisfy
(1)$ s_{tj} = f(R_t, j)$, and (2)
$ \sum_{j=1}^J N_j f(R_t, j) =0   .$
Condition (1) builds in consumer optimization; condition
(2) says that aggregate net savings equals zero (borrowing
equals lending).

 For the case in which the endowments, preferences, and group
sizes are constant across time, the interest rate is determined
at the intersection of the aggregate savings function
with the $R$ axis, depicted as $R_1$ in Figure \Fg{savings1f}. %Figure 8.7.
 No intergenerational   transfers occur in the nonmonetary
equilibrium.  The equilibrium consists of a sequence of separate
 two-period pure consumption loan economies of a type
analyzed by Irving Fisher (1907).
\auth{Fisher, Irving}

\subsection{Monetary equilibrium}

In an equilibrium with valued fiat currency, at each
date $t \geq 1$ the old receive goods from the young
in exchange for the currency stock $H$.
 For any variable $x$,  $\vec x = \{x_t\}_{t=1}^\infty$.
  An equilibrium with valued fiat money is a set of
sequences $\vec R, \vec p, \vec s$
such that (1) $\vec p$ is a positive sequence, (2) $R_t =p_t / p_{t+1}$,
(3) $s_{jt} = f( R_t,j)$,  and  (4)
$ \sum_{j=1}^J N_j f(R_t, j) = {H \over p_t}.$
Condition (1) states that currency is valued at all dates.
Condition (2) states  that currency and consumption loans
are perfect substitutes.  Condition (3) requires that saving
decisions are optimal. Condition (4) equates the
net saving of the young (the left side) to the net dissaving
of the old (the right side). The old supply currency inelastically.

We can determine a stationary equilibrium graphically.  A stationary
equilibrium satisfies $p_t =p$ for all $t$, which implies
$R =1$ for all $t$.  Thus, if it exists, a stationary equilibrium
solves
$$ \sum_{j=1}^J N_j f(1, j) = {H \over p}  \EQN equil4
 $$
for a positive price level. (See Figure \Fg{savings1f}.) %Figure 8.7.
Evidently, a stationary monetary equilibrium exists if the net
savings of the young are positive at $R=1$.

%%%%%%%%
%\topinsert{
%$$  \grafone{savings1.eps,height=2.5in}{{\bf Figure 8.7}
%The intersection of the aggregate savings function
%with a horizontal line at $R=1$ determines  a stationary equlibrium value of
%the price level, if positive.}
%$$
%}\endinsert
%%%%%%%%%%%%%%%

\midfigure{savings1f}
\centerline{\epsfxsize=3truein\epsffile{savings1.eps}}
\caption{The intersection of the aggregate savings function with a horizontal
line at $R=1$ determines  a stationary equilibrium value of the price level,
if positive.}
\infiglist{savings1f}
\endfigure

For the special case of logarithmic preferences and two
classes of young people, the aggregate savings function of the young is
time invariant and equal to
$$ \sum_j f(R,j) = .5(N_1 \alpha - N_2{\beta  \over R}). $$
Note that the equilibrium condition \Ep{equil4} can be written
$$ .5 N_1 \alpha = .5 N_2 {\beta  \over R} + {H  \over p}. $$
The left side is the demand for savings or the demand
for ``currency'' while the right side is the supply, consisting
of privately issued IOU's (the first term) and
government-issued currency.  The right side is thus
an abstract version of what is called  M1, which is a sum of privately
issued IOUs (demand deposits) and government-issued reserves and
currency.\index{M1}%

\subsection{Nonstationary equilibria}

Mathematically, the equilibrium conditions for the
 model with log preferences and
two groups have the same structure as the model analyzed
previously in equations \Ep{samdiff1} and \Ep{samsol1},
 with simple reinterpretations of
parameters.  We leave it to the reader here and in an exercise
to show that if there exists a stationary equilibrium with
valued fiat currency, then there exists a continuum of
equilibria with valued fiat currency, all but one of which
have the real value of government currency approaching
zero asymptotically.  A linear difference equation
like \Ep{samdiff1} supports this conclusion.

\subsection{The real bills doctrine}

In nineteenth-century Europe and the early days of the Federal Reserve system
in the United States, central banks conducted open market operations not by
purchasing government securities but by purchasing safe (risk-free)
short-term private IOUs.
 We now analyze this old-fashioned type of
open market operation. We allow the government to issue additional
currency each period.  It uses the proceeds exclusively to purchase
private IOUs (make loans to private agents) in the amount $L_t$ at time $t$.
 Such open market
operations are subject to the sequence of restrictions
$$ L_1 = {H_1 - H_0 \over p_1} $$
and 
$$L_t = R_{t-1} L_{t-1} + {H_t - H_{t-1} \over p_t} \EQN omo1  $$
for $t \geq 2$ and $H_0 = H>0$ given.
Here $L_t$ is the amount of the time $t$ consumption good that the government
lends to the private sector from period $t$ to period $t+1$.
Equation \Ep{omo1} states that the government finances these loans in two ways: first,
by rolling over the  proceeds
$R_{t-1} L_{t-1}$ from the repayment of last period's loans, and second, by injecting new  currency  in the
amount $H_t-H_{t-1}$.  With the government injecting new currency and purchasing loans in this
way each period, the equilibrium condition in the loan market becomes
$$ \sum_{j=1}^J N_j f(R_t, j) + L_t = {H_{t-1}\over p_t}
    + {H_{t}- H_{t-1} \over p_t}   \EQN equil5
$$
where the first term on the right is the real dissaving of the
old at $t$ (their real balances) and the second term
is the real value of the new money printed by the monetary
authority to finance purchases of private IOUs issued
by the young at $t$.  The left side is the net savings of
the young plus the savings of the government.
\index{open market operation!in private securities}
\index{real bills doctrine}
\auth{Smith, Adam}

 Under several guises, the effects of open market operations
like this have concerned monetary economists for centuries.\NFootnote{One
issue concerned the effects on the price
level of allowing banks to issue private bank notes.  Nothing in our model requires
that the notes $H_t$ be issued by the government.
We can also think of them as being issued by a private bank.}
We  state the following proposition:
\medskip
\noindent{\sc Irrelevance of Open Market Operations:}   
Positive sequences $\{L_t, H_t\}_{t=1}^\infty$
 that satisfy the constraint
\Ep{omo1} on open market operations  are associated with the same equilibrium allocation,
interest rate, and price level sequences.

\medskip
\noindent{\sc Proof:}
Evidently, we can write the equilibrium condition \Ep{equil5} as
$$ \sum_{j=1}^J N_j f(R_t, j) + L_t = {H_{t}\over p_t}. \EQN equil2 $$
Starting from $L_0 = 0$ and using $R_{t-1} = {p_{t-1} \over p_t}$
in equation \Ep{omo1} 
for $t \geq 2$ gives
$$\eqalign{ L_1 & = {H_1 - H_0 \over p_1} \cr 
            L_2 & = R_1 L_1 + {H_1 - H_0 \over p_1} = {H_2 - H_0\over p_2} }$$
%         
%iterating \Ep{omo1} once  and using $R_{t-1} = {p_{t-1} \over p_t}$
%gives
%$$ L_t = R_{t-1} R_{t-2} L_{t-2} + {H_t - H_{t-2} \over p_t} .$$
%Iterating back to time $0$ and using $L_0 = 0$ gives
and so on leading to
$$ L_t = {H_t - H_0\over p_t }. \EQN omo2 $$
Substituting \Ep{omo2} into \Ep{equil2} gives
$$ \sum_{j=1}^J N_j f(R_t, j)  = {H_{0}\over p_t}. \EQN equil3 $$
This is the same equilibrium condition in the economy with
no open market operations, i.e., the economy
with $L_t \equiv 0$ for all $t\geq 1$.   Any price level  and rate of
return sequence that solves \Ep{equil3} also solves \Ep{equil5} for
any $L_t$ sequence satisfying \Ep{omo1}. \qed

\medskip
This proposition captures  the spirit of Adam Smith's
real bills doctrine, which states that if the
government issues bank notes to purchase safe evidences
of private indebtedness, it is not inflationary. The quantity
$ {H_t - H_0\over p_t }$ in formula \Ep{omo2} is sometimes called ``inside money'' 
that is 100\% backed by evidences of private indebtedness while
$H_0$ is called ``outside money'' that is ``unbacked''.   
Sargent and Wallace (1982) analyze
settings in which the money market is initially separated from
the credit market by some legal restrictions that inhibit
intermediation.  Then open market operations are no longer
irrelevant because they can  partially  undo the
legal restrictions.  Sargent and Wallace show how those
legal restrictions can help stabilize the
price level at  costs in terms of economic efficiency.
Kahn and Roberds (1998) extend the Sargent and Wallace model  to study
 issues about regulating electronic payments systems.
\auth{Roberds, William}
\auth{Kahn, Charles}
\auth{Wallace, Neil}
\auth{Smith, Lones}
\auth{Kandori, Michihiro}
\section{Gift-giving equilibrium}
Michihiro Kandori (1992) and Lones Smith (1992) have used ideas
from the literature on reputation (see chapter \use{credible})
to study whether there exist history-dependent sequences of gifts that
support an optimal allocation. \index{gift-giving game!with overlapping
generations}
 Their idea is to set up the economy as a game played with a sequence
of players.  We briefly describe a gift-giving game for an overlapping
generations economy in which voluntary intergenerational gifts support
an optimal allocation.
Suppose that the consumption of an initial old person is
$$  c^0_1 = y_1^0 + s_1$$
 and the utility of each young agent is
$$ u(y^i_i - s_i) + u(y^i_{i+1} + s_{i+1}), \quad i \geq 1
  \EQN utility1 $$
where $s_i\geq 0$ is the gift from a young person
at $i$ to an old person at $i$.  Suppose that the endowment pattern is
$y^i_i = 1-\epsilon, y^i_{i+1} = \epsilon$, where
$\epsilon \in (0,.5)$.

Consider the following system of expectations, to which a young
person chooses whether to conform:
$$\EQNalign{ s_i& = \cases{ .5 - \epsilon & if $v_i = \overline v$; \cr
                             0 &otherwise.\cr } \EQN folk1;a  \cr
             v_{i+1}& = \cases{ \overline v & if $v_i = \overline v$ and
                                       $s_i = .5 - \epsilon$; \cr
                               \underline v & otherwise.\cr} \EQN folk1;b \cr}$$
Here we  are free  to take
$ \overline v = 2 u(.5)$ and $\underline v= u(1-\epsilon) +
u(\epsilon)$.    These are ``promised utilities.'' We make them serve
as ``state variables'' that summarize the history of intergenerational
gift giving.   To start, we need an initial value
$v_1$.   Equations   \Ep{folk1} act as the transition
laws that young agents face in choosing $s_i$ in
\Ep{utility1}.

  An initial condition $v_1$ and the rule \Ep{folk1} form a system
of expectations that tells the young person of each generation
what he is  expected to give.  His gift is  immediately handed over
to an old person.
A system of expectations is called an {\it equilibrium}
if for each $i \geq 1$, each young   agent chooses
to conform.

  We can immediately compute two equilibrium systems of
expectations.  The first  is
the ``autarky'' equilibrium: give nothing yourself and
expect  all future generations to give nothing. To  represent
this equilibrium within equations \Ep{folk1}, set $v_1 \neq \overline v$.
It is easy to verify that each young person will
confirm what is expected of him in this equilibrium.
Given that future generations will not give,  each young
person chooses not to give.

  For the second equilibrium, set $v_1 = \overline v$.   Here
each consumer chooses to give the expected amount, because
failure to do so causes  the next generation of
young people not to give; whereas affirming the expectation
to give passes that expectation along to the next generation,
which affirms it in turn.
Each of these equilibria is credible, in the sense
of subgame perfection, to be studied extensively in chapter
\use{credible}.


\auth{Kocherlakota, Narayana R.}
 Narayana Kocherlakota (1998) has  compared gift giving and monetary
equilibria in a variety of environments and has used the comparison to
provide a  precise sense in which ``money'' substitutes for ``memory''.

\auth{Mankiw, Gregory}
\auth{Summers, Lawrence}
\auth{Zeckhauser, Richard}
\auth{Abel, Andrew}
\auth{Diamond, Peter A.}
\section{Concluding remarks}
The overlapping generations model is a workhorse in
analyses of  public finance, welfare economics,
and demographics.  Diamond (1965) studied
some fiscal policy issues within a version of the model with a neoclassical production.
He showed that, depending on preference and productivity
parameters, equilibria of the model can
 have too much capital, and that such capital
overaccumulation can be corrected by having the government
issue and perpetually roll over unbacked debt.\NFootnote{Abel,
Mankiw, Summers, and Zeckhauser (1989) propose an empirical
test of whether there is capital overaccumulation in the U.S.
economy, and conclude that there is not.}
Auerbach  and  Kotlikoff (1987)  formulated a long-lived overlapping
generations model with capital, labor, production, and
various kinds of taxes. They used the model to study
a host of fiscal issues. Rios-Rull (1994a) used a calibrated
overlapping generations growth model to examine
the quantitative importance of market incompleteness for
insuring against aggregate risk. See Attanasio (2000) for
a  review of theories and evidence
about consumption within life-cycle models.
\auth{Attanasio, Orazio}

Several authors in a 1980 volume edited by John Kareken
and Neil Wallace argued through example that the overlapping
generations model is useful for analyzing a variety of issues in
monetary economics.  We refer to that volume, McCandless
and Wallace (1992), Champ and Freeman (1994), Brock (1990),
and Sargent (1987b) for a variety of applications
of the overlapping generations model to issues in monetary
economics.
\auth{Wallace, Neil}
\auth{Kareken, John}
\auth{Champ, Bruce}
\auth{Freeman, Scott}
\auth{Brock, William A.}
\auth{McCandless, George T.}

%\section{Exercises}
\showchaptIDfalse
\showsectIDfalse
\section{Exercises}
\showchaptIDtrue
\showsectIDtrue
\noindent{\it Exercise \the\chapternum.1} \ \  At each date $t \geq 1$, an  economy
consists of  overlapping generations of  a constant
number $N$ of two-period-lived agents.  Young
agents born in $t$ have preferences over consumption streams of
a single good that are ordered by
$ u(c_t^t) + u(c^t_{t+1}) $, where
$u(c) = c^{1 - \gamma} / (1-\gamma)$, and where $c^i_t$ is the
consumption of an agent born at $i$ in time $t$. It is understood
that $\gamma >0$, and that when $\gamma = 1$, $u(c) = \ln c$.   Each
young agent born at $t \geq 1$ has identical preferences and endowment
pattern
$(w_1,w_2)$, where $w_1$  is the endowment when young and  $w_2$
 is the endowment when old.
Assume $0 < w_2 < w_1$.
 In addition,
there are some initial old agents at time $1$ who
are endowed with $w_2$ of the time $1$ consumption good, and
who order consumption streams by $c^0_1$.   The initial old
(i.e., the old  at $t=1$)
are also endowed with $M$ units of unbacked fiat currency.  The
stock of currency is constant over time.
\medskip
\noindent{\bf a.}    Find the saving function of a young agent.
\medskip
\noindent{\bf b.}  Define an equilibrium with valued fiat
currency.
\medskip
\noindent{\bf c.}   Define a stationary equilibrium with valued
fiat currency.
\medskip
\noindent{\bf d.}  Compute a stationary equilibrium with valued
fiat currency.
\medskip
\noindent{\bf e.}  Describe how many equilibria with
valued fiat currency there are.  (You are not being asked to
compute them.)
\medskip
\noindent{\bf f.}  Compute the limiting value as $t \rightarrow + \infty$
of the rate of return on currency in each of the nonstationary
equilibria with valued fiat currency.  Justify your calculations.

\bigskip
\noindent{\it Exercise \the\chapternum.2} \ \
  Consider an economy with overlapping generations
of a constant population of an even number  $N$ of  two-period-lived agents.
New young agents are born at each date $t\geq 1$.
Half of the young agents are endowed with
$w_1$ when young and $0$ when old.  The other half are
endowed with $0$ when young and  $w_2$ when old.
Assume $0 < w_2 < w_1$.
Preferences of all young agents are as in problem
1, with $\gamma =1$.   Half of the $N$ initial old are endowed
with $w_2$  units of the consumption good and half are endowed
with nothing.  Each old  person orders consumption streams by
$c^0_1$.  Each old person at $t=1$ is endowed with $M$ units
of unbacked fiat currency. No other generation is endowed with
fiat currency.  The stock of fiat currency is fixed over time.

\medskip
\noindent{\bf a.}  Find the saving function of each of the two
types of young person for $t\geq 1$.
\medskip\noindent {\bf b.}  Define an equilibrium without valued fiat currency.
Compute  all such equilibria.
\medskip
\noindent{\bf c.}  Define an equilibrium with valued fiat currency.
\medskip
\noindent{\bf d.}  Compute all the (nonstochastic) equilibria
with valued fiat currency.
\medskip
\noindent{\bf e.}  Argue that there is a unique stationary equilibrium
with valued fiat currency.
\medskip
\noindent{\bf f.}  How are the various equilibria  with valued
fiat currency ranked by
the Pareto criterion?

\bigskip
\noindent{\it Exercise \the\chapternum.3} \ \   Take the economy of exercise
{\it \the\chapternum.1\/}, but make
one change.    Endow the initial old with a tree that
 yields a constant  dividend of $d >0$ units of the consumption
good for each $t \geq 1$.
\medskip
\noindent{\bf a.}   Compute all  the equilibria with valued
fiat currency.
\medskip
\noindent{\bf b.}  Compute all  the equilibria without valued
fiat currency.
\medskip
\noindent{\bf c.}  If you want,  you can  answer both parts
of this question in the context of the following
particular  numerical
example:  $w_1 =10,   w_2 =5, d=.000001$.


\bigskip
\noindent{\it Exercise  \the\chapternum.4} \ \  Take the economy of exercise {\it \the\chapternum.1\/} and make
the following two changes. First, assume
that $\gamma=1$. Second, assume that the number of young
agents born at $t$ is $N(t) =n N(t-1)$, where $N(0) >0$ is
given and $n \geq 1$.    Everything else about the economy remains
the same.

\medskip
\noindent{\bf a.}  Compute   an equilibrium without valued
fiat money.
\medskip
\noindent{\bf b.}  Compute a stationary equilibrium  with valued
fiat money.
\bigskip
\noindent{\it Exercise \the\chapternum.5} \ \  Consider an economy consisting of overlapping
generations of two-period-lived consumers.  At each date
$t\geq 1$ there are born $N(t)$ identical young people each of whom is
endowed with $w_1 >0$ units of a single consumption good when young
and $w_2 >0$ units of the consumption good when old.   Assume
that $w_2 < w_1$.    The consumption good is not storable.
The population of young people is described by
$N(t) = n N(t-1)$, where $n > 0$.   Young people
born at $t$ rank utility streams according to
$\ln(c^t_t) + \ln(c^t_{t+1})$ where $c^i_t$ is the consumption
of the time $t$ good of an agent born in $i$.    In addition,
there are $N(0)$ old people at time $1$, each of whom is
endowed with $w_2$ units of the time $1$ consumption good.
The old at $t=1$ are also endowed with one unit of  unbacked
pieces of infinitely durable but  intrinsically worthless pieces
of paper called fiat money.
\medskip
\noindent {\bf a.}  Define an equilibrium without valued fiat
currency.   Compute such an equilibrium.

\medskip
\noindent{\bf b.}  Define an equilibrium with valued fiat currency.
\medskip
\noindent{\bf c.}  Compute all equilibria with valued fiat currency.
\medskip
\noindent{\bf d.}  Find the limiting rates of return on currency
as $t \rightarrow + \infty$ in
each of the equilibria that you found in part   c.  Compare them
with the one-period  interest rate in the
equilibrium in part a.
\medskip
\noindent{\bf e.}  Are the equilibria in part c ranked according
to the Pareto criterion?


\medskip
\noindent{\it Exercise \the\chapternum.6} \ \  {\bf Exchange rate determinacy}
\medskip
\noindent   The world consists of two economies, named $i=1,2$, which
except for their governments' policies are ``copies'' of
one another.  At each date $t\geq 1$,
there is a single consumption good, which is storable, but
only  for rich people.  Each economy
consists of overlapping generations of two-period-lived agents.
At each $t \geq 1$, in economy $i$, $N$ poor people and
$N$ rich people are born.
Let $c_t^h(s),y_t^h(s) $
be the  time $s$ (consumption, endowment) of a type $h$ agent born
at $t$.  Poor agents are endowed with $[y_t^h(t), y_t^h(t+1)] = (\alpha,0)$;
rich agents are endowed
 with $[y_t^h(t), y_t^h(t+1)] = (\beta,0)$, where $\beta > > \alpha$.
In each country, there are $2N$ initial old who are
endowed in the aggregate
with $H_i(0)$ units of an unbacked currency and with
$2 N \epsilon$ units of the time $1$ consumption good.
For the rich people, storing $k$ units of the time $t$
 consumption good produces $R k$ units of the time $t+1$
 consumption good, where $R >1$ is a fixed gross
rate of return  on storage.  Rich people can earn
the rate of return $R$ either by storing goods or by lending to either
government by means of indexed bonds.
      We assume that poor people are prevented from storing
capital or holding indexed government debt by the sort
of denomination and intermediation restrictions
described by Sargent and Wallace (1982).

 At each $t\geq 1$, all young agents
order consumption streams according to
$\ln c_t^h(t) + \ln c_t^h(t+1)$.

   For $t \geq 1$, the government of country $i$ finances
a stream of purchases (to be thrown into the ocean)
of $G_i(t)$ subject to the following budget constraint:
$$ G_i(t) + R B_i(t-1) = B_i(t) + {H_i(t) - H_i(t-1)\over p_i(t)}
                        + T_i(t), \leqno(1)$$
where $B_i(0)=0$;
$p_i(t)$ is the price level in country $i$;
 $T_i(t)$ are lump-sum taxes levied by the government
on the {\it rich} young people at time $t$; $H_i(t)$ is
the stock of $i$'s fiat currency at the end
of period $t$; $B_i(t)$ is the stock of indexed government
interest-bearing debt (held by the rich of either country).   The government
does not explicitly tax poor people, but might
tax then through an inflation tax.  Each government levies
a lump-sum tax of $T_i(t)/N$ on each young rich citizen of its
own country.


  Poor people in both countries  are free to hold whichever
currency they prefer.  Rich people can hold debt of either
government and can also store; storage and both government debts
bear a  constant gross rate of return $R$.
\medskip

\noindent{\bf a.}  Define an {\it equilibrium\/} with valued fiat
currencies (in both countries).
\medskip
\noindent{\bf b.}  In a nonstochastic equilibrium, verify the following
proposition: if an equilibrium exists in which both fiat currencies
are valued, the exchange rate between the two currencies must be
constant over time.
\medskip


\noindent{\bf c.}  Suppose that government policy   in each
country is characterized by specified (exogenous)
levels $G_i(t) = G_i, T_i(t) = T_i$,
$B_i(t) = 0, \forall t \geq 1$. (The remaining elements of government policy
adjust to   satisfy the government budget constraints.)
Assume that the exogenous components of policy have
been set so that an equilibrium with two valued fiat
currencies exists.    Under this
description of policy, show that the equilibrium exchange rate
is indeterminate.

\medskip
\noindent{\bf d.}  Suppose that government policy in each country is
described as follows: $G_i(t)=G_i, T_i(t)=T_i, H_i(t+1)=H_i(1),
B_i(t)=B_i(1) \
 \forall t \geq 1$.
Show that if there exists an equilibrium with two valued fiat
currencies, the exchange rate is determinate.

\medskip

\noindent{\bf e.}  Suppose that government policy in country
$1$ is specified in terms of exogenous levels of
$s_1 = [H_1(t) - H_1(t-1)]/p_1(t)  \ \forall t \geq 2$,
and $G_1(t)=G_1 \ \forall t \geq 1$.  For country $2$,
government policy consists of exogenous levels
of $B_2(t)= B_2(1), G_2(t) = G_2 \forall t \geq 1$.  Show that if there
exists an equilibrium with two valued fiat currencies, then
the exchange rate is determinate.


\medskip

\noindent{\it Exercise \the\chapternum.7}\quad {\bf Credit controls}
\medskip
\noindent Consider the following overlapping generations model.  At each date $t\ge 1$
there appear $N$ two-period-lived young people, said to be of generation $t$,
who live and consume during periods $t$ and $(\toone)$.  At time $t=1$ there
exist $N$ old people who are endowed with $H(0)$ units of paper ``dollars,''
which they offer to supply inelastically to the young of generation 1 in
exchange for goods.  Let $p(t)$ be the price of the one good in the model,
measured in dollars per time $t$ good.  For each $t\ge 1$, $N/2$ members of
generation $t$ are endowed with $y>0$ units of the good at $t$ and 0 units at
$(\tone)$, whereas the remaining $N/2$ members of generation $t$ are endowed
with 0 units of the good at $t$ and $y>0$ units when they are old.  All members
of all generations have the same utility function:
$$u[c_t^h(t), c_t^h(\tone)]=\ln c_t^h(t) +\ln c_t^h(\tone),$$
where $c_t^h(s)$ is the consumption of agent $h$ of generation $t$ in period
$s$.  The old at  $t=1$ simply maximize $c_0^h(1)$.  The consumption good is
nonstorable.  The currency supply is constant through time, so $H(t)=H(0)$,
$t\ge 1$.
\medskip
\noindent{\bf a.} Define a competitive equilibrium without valued currency for this
model.  Who trades what with whom?
\medskip
\noindent{\bf b.} In the equilibrium without valued fiat currency,
 compute   competitive equilibrium values of the
gross return on consumption loans, the consumption allocation of the old at
$t=1$, and that of the ``borrowers'' and ``lenders'' for $t\ge 1$.
\medskip
\noindent{\bf c.} Define a competitive equilibrium with valued currency.  Who trades
what with whom?\medskip
\noindent{\bf d.} Prove that for this economy there does not exist a competitive
equilibrium with valued currency.\medskip
\noindent{\bf e.} Now suppose that the government imposes the restriction that
$l_t^h(t) [1+r(t)]\ge -y/4$, where $l_t^h(t)[1+r(t)]$ represents claims on
$(\tone)$--period consumption purchased (if positive) or sold (if negative) by
consumer $h$ of generation $t$.  This is a restriction on the amount of
borrowing.  For an equilibrium without valued currency, compute the consumption
allocation and the gross rate of return on consumption loans.
\medskip
\noindent{\bf f.} In the setup of part  e,
 show that there exists an equilibrium with
valued currency in which the price level obeys the quantity theory equation
$p(t)=qH(0)/N$.  Find a formula for the undetermined coefficient $q$.  Compute
the consumption allocation and the equilibrium rate of return on consumption
loans.\medskip
\noindent{\bf g.} Are lenders better off in economy b or economy f?  What about
borrowers?  What about the old of period 1 (generation 0)?
\medskip
\noindent{\it Exercise \the\chapternum.8}\quad {\bf Inside money and real bills}
\medskip

\noindent Consider the following overlapping generations model of two-period-lived
people.  At each date $t\ge 1$ there are born $N_1$ individuals of type 1 who
are endowed with $y>0$ units of the consumption good when they are young and
zero units when they are old; there are also born $N_2$ individuals of type 2
who are endowed with zero units of the consumption good when they are young and
$Y>0$ units when they are old.  The consumption good is nonstorable.  At time
$t=1$, there are $N$ old people, all of the same type, each endowed with zero
units of the consumption good and $H_0/N$ units of unbacked paper called ``fiat
currency.'' The populations of type 1 and 2 individuals, $N_1$ and $N_2$,
remain constant for all $t\ge 1$.  The young of each generation are identical
in preferences and maximize the utility function $\ln c_t^h(t) +\ln
c_t^h(\tone)$ where $c_t^h(s)$ is consumption in the $s$th period of a member
$h$ of generation $t$.
\medskip
\noindent{\bf a.} Consider the equilibrium without valued currency (that is, the
equilibrium in which there is no trade between generations).  Let $[1+r(t)]$ be
the gross rate of return on consumption loans.  Find a formula for $[1+r(t)]$
as a function of $N_1, N_2,y$, and $Y$.\medskip
\noindent{\bf b.} Suppose that $N_1,N_2,y$,
 and $Y$ are such that $[1+r(t)]>1$ in the
equilibrium without valued currency.  Then prove that there can
exist no quantity-theory-style equilibrium where fiat currency is
valued and where the price level $p(t)$ obeys the quantity theory
equation $p(t)=q\cdot H_0$, where $q$ is a positive constant and
$p(t)$ is measured in units of currency per unit good.
\medskip
\noindent{\bf c.} Suppose that $N_1,N_2,y$, and $Y$ are such that in the
nonvalued-currency equilibrium $1+r(t)<1$.  Prove that there
exists an equilibrium in which fiat currency is valued and that
there obtains the quantity theory equation $p(t)=q\cdot H_0$,
where $q$ is a constant.  Construct an argument to show that the
equilibrium with valued currency is not Pareto superior to the
nonvalued-currency equilibrium.
\medskip \noindent{\bf d.} Suppose that
$N_1,N_2,y$, and $Y$ are such that, in the preceding
nonvalued-currency economy, $[1+r(t)]<1$, there exists an
equilibrium in  which fiat currency is valued.  Let $\bar p$ be
the stationary equilibrium price level in that economy.  Now
consider an alternative economy, identical with the preceding one
in all respects except for the following feature: a government
each period purchases a constant amount $L_g$ of consumption loans
and pays for them by issuing debt on itself, called ``inside
money'' $M_I$, in the amount $M_I(t)=L_g\cdot p(t)$.  The
government never retires the inside money, using the proceeds of
the loans to finance new purchases of consumption loans in
subsequent periods.  The quantity of outside money, or currency,
remains $H_0$, whereas the ``total high-power money'' is now
$H_0+M_I(t)$. \item{(i)} Show that in this economy there
exists a valued-currency equilibrium in which the price level is
constant over time at $p(t)=\bar p$, or equivalently, with $\bar
p=qH_0$ where $q$ is defined in part c. \item{(ii)} Explain
why government purchases of private debt are not inflationary in
this economy. \item{(iii)} In many models,
once-and-for-all government open-market operations in private debt
normally affect real variables and/or price level.  What accounts
for the difference between those models and the one in this
exercise?
\newpage
\medskip
\noindent{\it Exercise \the\chapternum.9}\quad {\bf Social security and the price level}
\medskip
\noindent Consider an economy (``economy I'') that consists of overlapping generations of
two-period-lived people.  At each date $t\ge 1$ there is born a constant
number $N$ of young people, who desire to consume both when they are young, at
$t$, and when they are old, at $(\tone)$.  Each young person has the utility
function $\ln c_t(t)+\ln c_t(\tone)$, where $c_s(t)$ is time $t$ consumption of
an agent born at $s$.  For all dates $t\ge 1$, young people are endowed with
$y>0$ units of a single nonstorable consumption good when they are young and
zero units when they are old.  In addition, at time $t=1$ there are $N$ old
people endowed in the aggregate with $H$ units of unbacked fiat currency.  Let
$p(t)$ be the nominal price level at $t$, denominated in dollars per time $t$
good.
\medskip\noindent{\bf a.}
 Define and compute an equilibrium with valued fiat currency for this
economy.  Argue that it exists and is unique.  Now consider a second economy
(``economy II'') that is identical to economy I except that economy II
possesses a social security system.  In particular, at each date $t\ge 1$, the
government taxes $\tau>0$ units of the time $t$ consumption good away from each
young person and at the same time gives $\tau$ units of the time $t$
consumption good to each old person then alive.
\medskip\noindent{\bf b.} Does economy II possess an equilibrium with valued fiat currency?
Describe the restrictions on the parameter $\tau$, if any, that are needed to
ensure the existence of such an equilibrium.
\medskip\noindent{\bf c.} If an equilibrium with valued fiat currency exists, is it unique?
\medskip\noindent{\bf d.}
 Consider the {\it stationary\/} equilibrium with valued fiat currency.
Is it unique? Describe how the value of currency or price level would vary
across economies with differences in the size of the social security system, as
measured by $\tau$.
\medskip
\noindent{\it Exercise \the\chapternum.10}\quad {\bf Seigniorage}
\medskip
\noindent Consider an economy consisting of overlapping generations of two-period-lived
agents.  At each date $t\ge 1$, there are born $N_1$ ``lenders'' who are
endowed with $\a>0$ units of the single consumption good when they are young
and zero units when they are old.  At each date $t\ge 1$, there are also born
$N_2$ ``borrowers'' who are endowed with zero units of the consumption good
when they are young and $\be>0$ units when they are old.  The good is
nonstorable, and $N_1$ and $N_2$ are constant through time.  The economy starts
at time 1, at which time there are $N$ old people who are in the aggregate
endowed with $H(0)$ units of unbacked, intrinsically worthless pieces of paper
called dollars.  Assume that $\a,\be,N_1$, and $N_2$ are such that
$${N_2\be\over N_1\a} <1.$$
Assume that everyone has preferences
$$u[c_t^h(t), c_t^h(\tone)]=\ln c_t^h(t) +\ln c_t^h(\tone),$$
where $c_t^h(s)$ is consumption of time $s$ good of agent $h$ born at time $t$.
\medskip\noindent{\bf a.} Compute the equilibrium interest rate on consumption loans in the
equilibrium without valued currency.
\medskip\noindent{\bf b.} Construct a {\it brief\/} argument to establish whether or not the
equilibrium without valued currency is Pareto optimal.

The economy also contains a government that purchases and destroys $G_t$ units
of the good in period $t$, $t\ge 1$.  The government finances its purchases
entirely by currency creation.  That is, at time $t$,
$$G_t ={H(t)-H(t-1)\over p(t)},$$
where $[H(t)-H(t-1)]$ is the additional dollars printed by the government at
$t$ and $p(t)$ is the price level at $t$.  The government is assumed to
increase $H(t)$ according to
$$H(t)=zH(t-1),\qquad z\ge 1,$$
where $z$ is a constant for all time $t\ge 1$.

At time $t$, old people who carried $H(t-1)$ dollars between $(t-1)$ and
$t$ offer these $H(t-1)$ dollars in exchange for time $t$ goods.  Also at $t$
the government offers $H(t)-H(t-1)$ dollars for goods, so that $H(t)$ is the
total supply of dollars at time $t$, to be carried over by the young into time
$(\tone)$.
\medskip\noindent{\bf c.} Assume that $1/z>N_2\be/N_1\a$.  Show that under this assumption
there exists a continuum of equilibria with valued currency.
\medskip\noindent{\bf d.}
 Display the unique stationary equilibrium with valued currency in the
form of a ``quantity theory'' equation.  Compute the equilibrium rate of return
on currency and consumption loans.
\medskip\noindent{\bf e.} Argue that if $1/z<N_2\be/N_1\a$, then there exists no
valued-currency equilibrium.  Interpret this result.  ({\it Hint:}
 Look at the rate
of return on consumption loans in the equilibrium without valued currency.)
\medskip\noindent{\bf f.}
 Find the value of $z$ that {\it maximizes\/} the government's $G_t$ in
a stationary equilibrium.  Compare this with the largest value of $z$ that is
compatible with the existence of a valued-currency equilibrium.
\medskip
\noindent{\it Exercise \the\chapternum.11} \quad {\bf Unpleasant monetarist arithmetic}
\medskip
\noindent Consider an economy in which the aggregate demand for government currency for
$t\ge 1$ is given by $[M(t)p(t)]^d =g[R_1(t)]$, where $R_1(t)$ is the gross
rate of return on currency between $t$ and $(\tone)$, $M(t)$ is the stock of
currency at $t$, and $p(t)$ is the value of currency in terms of goods at $t$
(the reciprocal of the price level).  The function $g(R)$ satisfies
$$g(R) (1-R)  =h(R)>0\qquad {\rm for}\ \ R\in (\underline{R},1),
\leqno(1) $$
where $h(R)  \le 0$  for $R<\underline{R}, R\ge 1,
\underline{R}>0$ and
where
$h'(R)<0$ for $R>R_m$, $ h'(R)>0$ for $R<R_m$
$h(R_m)>D$, where $D$ is a positive number to be defined
shortly.
The government faces an infinitely elastic demand for its interest-bearing
bonds at a constant-over-time gross rate of return $R_2>1$.  The government
finances a budget deficit $D$, defined as government purchases minus explicit
taxes, that is constant over time.  The government's budget constraint is
$$D=p(t) [M(t)-M(t-1)] +B(t) -B(t-1) R_2,\qquad t\ge 1,\leqno(2)$$
subject to $B(0)=0, M(0)>0$.  In equilibrium,
$$M(t) p(t)=g[R_1(t)].\leqno(3)$$
The government is free to choose paths of $M(t)$ and $B(t)$, subject to
equations (2) and (3).
\medskip\noindent{\bf a.}
 Prove that, for $B(t)=0$, for all $t>0$, there exist two stationary
equilibria for this model.
\medskip\noindent{\bf b.}
 Show that there exist values of $B>0$, such that there exist
stationary equilibria with $B(t)=B$, $M(t)p(t)=Mp$.
\medskip\noindent{\bf c.} Prove a version of the following proposition: among stationary
equilibria, the lower the value of $B$, the lower the stationary rate of
inflation consistent with equilibrium. (You will have to make an assumption
about Laffer curve effects to obtain such a proposition.)

This problem displays some of the ideas used by Sargent and Wallace (1981).
They argue that, under assumptions like those leading to the
proposition stated in part c, the ``looser'' money is today [that is, the
higher $M(1)$ and the lower $B(1)$], the lower the stationary inflation rate.
\medskip
\noindent{\it Exercise \the\chapternum.12}\quad {\bf Grandmont-Hall}
\medskip
\noindent Consider a nonstochastic, one-good overlapping generations model consisting of
two-period-lived young people born in each $t\ge 1$ and  an initial group of
old people at $t=1$ who are endowed with $H(0)>0$ units of unbacked currency at
the beginning of period 1.  The one good in the model is not storable.  Let the
aggregate first-period saving function of the young be time-invariant and be
denoted $f[1+r(t)]$ where $[1+r(t)]$ is the gross rate of return on consumption
loans between $t$ and $(\tone)$.  The saving function is assumed to satisfy
$f(0)=-\infty$, $f'(1+r)>0$, $f(1)>0$.

Let the government pay interest on currency, starting in period 2 (to holders
of currency between periods 1 and 2).  The government pays interest on currency
at a nominal rate of $[1+r(t)] p(\tone)/\bar p$, where $[1+r(t)]$ is the real
gross rate of return on consumption loans, $p(t)$ is the price level at $t$,
and $\bar p$ is a target price level chosen to satisfy
$$\bar p=H(0)/f(1).$$
The government finances its interest payments by printing new money, so that
the government's budget constraint is
$$H(\tone)-H(t) =\left\{ [1+r(t)] {p(\tone)\over \bar p} -1\right\} H(t),\qquad
t\ge 1,$$
given $H(1)=H(0)>0$.  The gross rate of return on consumption loans in this
economy is $1+r(t)$.  In equilibrium, $[1+r(t)]$ must be at least
as great as the real rate of return on currency
$$1+r(t) \geq [1+r(t)] p(t)/\bar p =[1+r(t)] {p(\tone)\over \bar p} {p(t)\over p(\tone)}$$
with equality if currency is valued,
$$1+r(t) = [1+r(t)]p(t)/\bar p,\qquad 0<p(t)<\infty.$$
The loan market-clearing condition in this economy is
$$f[1+r(t)]=H(t)/p(t).$$
\medskip\noindent{\bf a.} Define an equilibrium.
\medskip\noindent{\bf b.}
 Prove that there exists a unique monetary equilibrium in this economy
and compute it.
\medskip
\noindent{\it Exercise \the\chapternum.13}\quad {\bf Bryant-Keynes-Wallace}
\medskip
\noindent Consider an economy consisting of overlapping generations of two-period-lived
agents.  There is a constant population of $N$ young agents born at each date
$t\ge 1$.  There is a single consumption good that is not storable.  Each agent
born in $t\ge 1$ is endowed with $w_1$ units of the consumption good when young
and with $w_2$ units when old, where $0<w_2<w_1$.  Each agent born at $t\ge 1$
has identical preferences $\ln c_t^h(t) +\ln c_t^h (\tone)$, where $c_t^h(s)$
is time $s$ consumption of agent $h$ born at time $t$.  In addition, at time 1,
there are alive $N$ old people who are endowed with $H(0)$ units of unbacked
paper currency and who want to maximize their consumption of the time 1 good.

A government attempts to finance a constant level of government purchases
$G(t)=G>0$ for $t\ge 1$ by printing new base money.  The government's budget
constraint is
$$G=[H(t)-H(t-1)]/p(t),$$
where $p(t)$ is the price level at $t$, and $H(t)$ is the stock of currency
carried over from $t$ to $(\tone)$ by agents born in $t$.  Let $g=G/N$ be
government purchases per young person.  Assume that purchases $G(t)$ yield no
utility to private agents.
\medskip\noindent{\bf a.} Define a stationary equilibrium with valued fiat currency.
\medskip\noindent{\bf b.} Prove that, for $g$ sufficiently small, there exists a stationary
equilibrium with valued fiat currency.
\medskip\noindent{\bf c.} Prove that, in general, if there exists one stationary equilibrium
with valued fiat currency, with rate of return on currency $1+r(t)=1+r_1$, then
there exists at least one other stationary equilibrium with valued currency
with $1+r(t)=1+r_2 \not=1+r_1$.
\medskip\noindent{\bf d.} Tell whether the equilibria described in parts  b and c
 are Pareto
optimal, among allocations among private agents of what is
left after the government takes $G(t)=G$ each period. (A proof is not required
here: an informal argument will suffice.)

Now let the government institute a forced saving program of the following form.
At time 1, the government redeems the outstanding stock of currency $H(0)$,
exchanging it for government bonds.  For $t\ge 1$, the government offers each
young consumer the option of saving at least $F$ worth of time $t$ goods in the
form of bonds bearing a constant rate of return $(1+r_2)$.  A legal prohibition
against private intermediation is instituted that prevents two or more private
agents from sharing one of these bonds.  The government's budget constraint for
$t\ge 2$ is
$$G/N =B(t) -B(t-1)(1+r_2),$$
where $B(t) \ge F$.  Here $B(t)$ is the saving of a young agent at $t$.  At
time $t=1$, the government's budget constraint is
$$G/N =B(1) -{H(0)\over Np (1)},$$
where $p(1)$ is the price level at which the initial currency stock is redeemed
at $t=1$.  The government sets $F$ and $r_2$.

Consider stationary equilibria with $B(t)=B$ for $t\ge 1$ and $r_2$ and $F$
constant.

\medskip\noindent{\bf e.}
 Prove that if $g$ is small enough for an equilibrium of the type
described in part a to
exist, then a stationary equilibrium with forced saving exists. (Either a
graphical argument or an algebraic argument is sufficient.)
\medskip\noindent{\bf f.} Given $g$, find the values of $F$ and $r_2$ that maximize the utility
of a representative young agent for $t\ge 1$.
\medskip\noindent{\bf g.} Is the equilibrium
 allocation associated with the values of $F$ and
$(1+r_2)$ found in part f optimal among those allocations that give $G(t)=G$ to
the government for all $t\ge 1$?  (Here an informal argument will suffice.)


\medskip
\noindent{\it Exercise \the\chapternum.14}\quad {\bf OLG, 1}
\medskip
\noindent An economy consists of overlapping generations of two-period live people. At dates $t \geq 1$ there are born
$N_t >0$ identical people who are endowed with $w_1 >0$ units of a nonstorable consumption good when they
are young and $w_2 >0 $ units of a nonstorable consumption good when they are old.  In addition, at time $1$ there
are $N_0> 0$ old people each of whom is endowed with ${M \over N_0}$ units of an intrinsically worthless unbacked (fiat) currency called
dollars (it is not in anyone's utility function and is useless for production).  
Population evolves as $N_{t+1} = n N_t$ for $t \geq 0$ where $n >0$.  The consumption profile
of a person born at time $t$ is $c_t^t$ when young and $c_{t+1}^t$ when old; he/she orders profiles according to the
utility functional
$ \ln (c_t^t) + \ln (c^t_{t+1}) .$

\medskip
\noindent 
{\bf a.} Define a competitive equilibrium in which no value is attached to fiat currency, being careful to define all
of the objects that comprise a competitive equilibrium and telling who trades what with whom.

\medskip
\noindent {\bf b.} Please compute all competitive equilibria in which no value is attached to fiat currency.

\medskip
\noindent{\bf c.} Please define a competitive equilibrium in which value is attached to fiat currency, 
being careful to define all
of the objects that comprise a competitive equilibrium and telling who trades what with whom.

\medskip
\noindent{\bf d.} A ``friend'' makes the following guess.  If $n w_1 > w_2$ there exist competitive equilibria
with valued fiat currency, while if $n w_1 < w_2$ there do not.  Please either confirm or dispose of this guess.
If the guess is true, please define and compute a {\it stationary\/} equilibrium with valued fiat currency and give formulas
for the equilibrium price level and rate of return on fiat currency.

\medskip
\noindent{\bf e.} Please argue that whenever a stationary equilibrium with valued fiat currency exists, then there
also exists a continuum of equilibria with valued fiat currency, in each of which the rate of return on currency
varies over time.  Please provide explicit formulas for an equilibrium  price level as a function of time.  
{\bf Hint:} It might help to deploy a linear difference equation.

\medskip
\noindent{\bf f.} Assuming that a nonstationary equilibrium with valued fiat currency exists, please compute the
the limiting value of the rate of return on currency as $t \rightarrow + \infty$. 




\medskip
\noindent{\it Exercise \the\chapternum.15}\quad {\bf OLG, 2}
\medskip
\noindent An economy consists of overlapping generations of two-period lived people.  Let $T > 2$ be a finite 
integer. At dates $1 \leq  t \leq T$ there are born
$N$ identical people who are endowed with $w_1 >0$ units of a nonstorable consumption good when they
are young and $w_2>0 $ units of a nonstorable consumption good when they are old.  
At dates $t \geq T+1$ there are born
$N$ identical people who are endowed with $\bar w_1 >0$ units of a nonstorable consumption good when they
are young and $\bar w_2>0$ units of a nonstorable consumption good when they are old. In addition, at time $1$ there
are $N> 0$ old people each of whom is endowed with ${M \over N}$ units of an intrinsically worthless unbacked (fiat) currency called
dollars (it is not in anyone's utility function and is useless for production). The consumption profile
of a person born at time $t$ is $c_t^t$ when young and $c_{t+1}^t$ when old; he/she orders profiles according to the
utility functional
$ \ln (c_t^t) + \ln (c^t_{t+1}) .$


\medskip
\noindent 
{\bf a.} Define a competitive equilibrium in which no value is attached to fiat currency, being careful to define all
of the objects that comprise a competitive equilibrium and telling who trades what with whom.

\medskip
\noindent {\bf b.} Please compute all competitive equilibria in which no value is attached to fiat currency.

\medskip
\noindent{\bf c.} Please define a competitive equilibrium in which value is attached to fiat currency, 
being careful to define all
of the objects that comprise a competitive equilibrium and telling who trades what with whom.

\medskip
\noindent {\bf d.} Please state conditions on $w_1, w_2, \bar w_1, \bar w_2$ for which there exist competitive
equilibria with valued fiat currency.

\medskip
\noindent {\bf e.} Under sufficient conditions for there to exist competitive equilibria with valued fiat currency,
please find explicit formulas for the equilibrium price level sequence in the equilibrium for which the value
of currency is highest for all $t \geq 1$. {\bf Hint:} It might help to deploy a linear difference equation.



\medskip
\noindent{\it Exercise \the\chapternum.16}\quad {\bf OLG, 3}
\medskip
\noindent   An economy consists of overlapping generations of two-period live people. At dates $t \geq 0$ there are born
$N_t >0$ identical people who are endowed with $w_1 >0$ units of a nonstorable consumption good when they
are young and $w_2 >0 $ units of a nonstorable consumption good when they are old.  Population evolves as $N_{t+1} = n N_t$ for $t \geq 0$ where $n >0$ and $N_0 >0$.  The consumption profile
of a person born at time $t$ is $c_t^t$ when young and $c_{t+1}^t$ when old.

\medskip
\noindent{\bf a.} Please show that for $t \geq 1$, the frontier of feasible allocations is described by
$$ c_t^{t-1} \le n(w_1- c_t^t) + w_2, \ c_t^{t-1} \geq 0, \ c_t^t \geq 0 .$$

\medskip
\noindent{\bf b.}  For $t \geq 1$, each young person  orders consumption profiles according to the
utility functional
$ \ln (c_t^t) + \ln (c^t_{t+1}) .$
Please derive the person's offer curve.

\medskip
\noindent {\bf c.} After defining a competitive equilibrium for this economy, please draw David Gale's diagram and use it to find a Pareto optimal competitive equilibrium  allocation.

\medskip
\noindent{\bf d.} Tell how to use the same diagram to find a competitive equilibrium allocation that is not Pareto optimal.


